\[
d ^ {+} (\mu , \gamma , \nu) = \dim \mathrm{E} ^ {+} (\mu , \gamma , \nu), \quad d ^ {-} (\mu , \gamma , \nu) = \dim \mathrm{E} ^ {-} (\mu , \gamma , \nu).
\]

For all \(\lambda \in \mathfrak{h}^*\) , put \(\lambda^{*} = -w_{0}\lambda\) , where \(w_{0}\) is the element of W that takes B to \(-B\) . Show that

\[
d ^ {+} (\mu , \gamma , \nu) = d ^ {-} (\mu , - \gamma^ {*}, \nu^ {*}).
\]

\begin{enumerate}
\def\labelenumi{\alph{enumi})}
\setcounter{enumi}{1}
\item
  Let \(\lambda_1, \lambda_2 \in \mathrm{P}_{++}\) , V the \(\mathfrak{g}\) -module \(\mathrm{Hom}_k(\mathrm{E}(\lambda_1^*), \mathrm{E}(\lambda_2))\) , U the set of \(\varphi \in V\) such that \(Y_\alpha. \varphi = 0\) for all \(\alpha \in B\) , and \(w\) a primitive vector in \(\mathrm{E}(\lambda_1^*)\) . Show that \(\varphi \mapsto \varphi(w)\) is an isomorphism from U to the set of \(v \in \mathrm{E}(\lambda_2)\) such that \(Y_\alpha^{\lambda_1^*(H_\alpha)+1}. v = 0\) for all \(\alpha \in B\) . (Use Exerc. 15 of §7 to prove surjectivity.)
\item
  With the notations of Prop. 2, prove that, for \(\lambda_1, \lambda_2, \lambda \in \mathrm{P}_{++}\) ,
\end{enumerate}

\[
\begin{array}{r} m (\lambda_ {1}, \lambda_ {2}, \lambda) = d ^ {+} (\lambda , \lambda_ {2} - \lambda_ {1} ^ {*}, \lambda_ {1} ^ {*}) = d ^ {-} (\lambda , \lambda_ {1} - \lambda_ {2} ^ {*}, \lambda_ {1}) \\ = d ^ {+} (\lambda_ {1}, \lambda - \lambda_ {2}, \lambda_ {2}) = d ^ {-} (\lambda_ {1}, \lambda_ {2} ^ {*} - \lambda^ {*}, \lambda_ {2} ^ {*}). \end{array}
\]

(Observe that \(m(\lambda_1,\lambda_2,\lambda)\) is the dimension of the space of \(\mathfrak{g}\) -invariant elements of

\[
\mathrm{E} (\lambda) ^ {*} \otimes \mathrm{E} (\lambda_ {1}) \otimes \mathrm{E} (\lambda_ {2}),
\]

and hence is equal to \(m(\lambda_1^*, \lambda, \lambda_2)\) .)

\(d)\) Let \(\lambda_1, \lambda_2 \in \mathrm{P}_{++}\) , and let \(\lambda\) be the unique element of \(\mathrm{P}_{++} \cap \mathrm{W}.\) ( \(\lambda_1 - \lambda_2^*\) ). We have

\[
m (\lambda_ {1}, \lambda_ {2}, \lambda) = 1.
\]

(Use c.) Deduce that \(\mathrm{E}(\lambda_1) \otimes \mathrm{E}(\lambda_2)\) contains a unique submodule isomorphic to \(\mathrm{E}(\lambda)\) .

\begin{enumerate}
\def\labelenumi{\alph{enumi})}
\setcounter{enumi}{4}
\tightlist
\item
  We retain the notations of \(d\) ). Show the equivalence of the following conditions:
\end{enumerate}

\begin{enumerate}
\def\labelenumi{(\roman{enumi})}
\item
  \(\lambda = \lambda_{1} + \lambda_{2}\)
\item
  \(\| \lambda \| = \| \lambda_1 + \lambda_2 \|\)
\item
  \(\lambda_{1}\) and \(\lambda_2^*\) are orthogonal
\item
  \(\lambda_{1}\) and \(\lambda_2^*\) are disjoint (Exerc. 13)
\item
  \(\lambda_{1}\) and \(\lambda_{2}\) are disjoint.
\end{enumerate}

Conclude that \(\mathrm{E}(\lambda_{1})\otimes\mathrm{E}(\lambda_{2})\) is not a simple module unless \(\lambda_{1}\) and \(\lambda_{2}\) are disjoint (and hence another proof of Exerc. 18 f) of §7).

\begin{enumerate}
\def\labelenumi{\arabic{enumi})}
\setcounter{enumi}{14}
\tightlist
\item
  Put \(\mathrm{N} = \mathrm{Card}(\mathrm{R}_{+})\) , \(c = \prod_{\alpha \in \mathrm{R}_{+}} \langle \rho, H_{\alpha} \rangle\) , and \(d(\lambda) = \dim \mathrm{E}(\lambda)\) if \(\lambda \in \mathrm{P}_{++}\) . Show that, for any real number \(s > 0\) ,
\end{enumerate}

\[
\sum_ {\lambda \in \mathrm{P} _ {+ +}} d (\lambda) ^ {- s} \leq \frac {1}{c} \left(\sum_ {m = 1} ^ {\infty} m ^ {- s}\right) ^ {\mathrm{N}}.
\]

Deduce that \(\sum_{\lambda \in \mathrm{P}_{+ + }}d(\lambda)^{-s} <   + \infty\) if \(s > 1\)

¶ 16) a) Take g to be of type \(F_{4}\) and use the notations of Chap. VI, Plate VIII. If i = 1, 2, 3, 4, put

\[
\mathcal {X} _ {i} = \mathrm{P} _ {+ +} \cap (\varpi_ {i} - \mathrm{Q} _ {+}).
\]

The set of weights of \(\mathrm{E}(\varpi_i)\) is the disjoint union of the \(\mathrm{W}\omega\) , where \(\omega\) belongs to \(\mathcal{X}_i\) (cf.~§7, Prop. 5 (iv)). We have:

\[
\mathcal {X} _ {1} = \{0, \varpi_ {1}, \varpi_ {4} \};
\]

\[
\mathcal {X} _ {2} = \{0, \varpi_ {1}, \varpi_ {2}, \varpi_ {3}, \varpi_ {4}, \varpi_ {1} + \varpi_ {4}, 2 \varpi_ {4} \};
\]

\[
\mathcal {X} _ {3} = \{0, \varpi_ {1}, \varpi_ {3}, \varpi_ {4} \};
\]

\[
\mathcal {X} _ {4} = \{0, \varpi_ {4} \}.
\]

\begin{enumerate}
\def\labelenumi{\alph{enumi})}
\setcounter{enumi}{1}
\tightlist
\item
  Show, by means of Th. 2, that:
\end{enumerate}

\(\dim \mathrm{E}(\varpi_1) = 52,\quad \dim \mathrm{E}(\varpi_2) = 1274,\quad \dim \mathrm{E}(\varpi_3) = 273,\quad \dim \mathrm{E}(\varpi_4) = 26.\)

\begin{enumerate}
\def\labelenumi{\alph{enumi})}
\setcounter{enumi}{2}
\tightlist
\item
  By using Chap. V, §3, Prop. 1, and the plates of Chap. VI, show that
\end{enumerate}

\[
\mathrm{Card} (\mathrm{W} \varpi_ {1}) = 2 ^ {7} 3 ^ {2} 2 ^ {- 3} (3!) ^ {- 1} = 24.
\]

Calculate \(\mathrm{Card}(\mathrm{W}\varpi_{2}),\ldots,\mathrm{Card}(\mathrm{W}.2\varpi_{4})\) similarly. Deduce that the number of weights of \(\mathrm{E}(\varpi_{2})\) is 553; since this number is strictly less than \(\dim\mathrm{E}(\varpi_{2})\) , one of these weights is of multiplicity \(\geq2\) .

\begin{enumerate}
\def\labelenumi{\alph{enumi})}
\setcounter{enumi}{3}
\item
  Make analogous calculations for \(\varpi_{1},\varpi_{3},\varpi_{4}\) . By using Exerc. 21 of §7, deduce that, if \(\rho\) is a non-zero simple representation of g, \(\rho\) admits a weight of multiplicity \(\geq2\) .
\item
  Prove the same result for a simple algebra of type \(\mathrm{E}_8\) .
\end{enumerate}

\(f)\) Let \(\mathfrak{a}\) be a splittable simple Lie algebra. Deduce from \(d\) , \(e)\) and Prop. 7 and 8 of §7 the equivalence of the following properties:

\begin{enumerate}
\def\labelenumi{(\roman{enumi})}
\item
  \(\mathfrak{a}\) admits a non-zero simple representation all of whose weights are of multiplicity 1;
\item
  \(\mathfrak{a}\) is neither of type \(\mathrm{F}_4\) nor of type \(\mathrm{E}_8\) .
\end{enumerate}

\exercisesection{\S~10}

\begin{enumerate}
\def\labelenumi{\arabic{enumi})}
\item
  Let \(\mathfrak{s} = \mathfrak{sl}(2, k)\) and \(\mathfrak{g} = \mathfrak{sl}(3, k)\) . Identify \(\mathfrak{s}\) with a subalgebra of \(\mathfrak{g}\) by means of an irreducible representation of \(\mathfrak{s}\) of degree 3. Show that every subspace of \(\mathfrak{g}\) containing \(\mathfrak{s}\) and stable under \(\operatorname{ad}_{\mathfrak{g}}\mathfrak{s}\) is equal to either \(\mathfrak{s}\) or \(\mathfrak{g}\) ; deduce that \(\mathfrak{s}\) is maximal among the subalgebras of \(\mathfrak{g}\) distinct from \(\mathfrak{g}\) .
\item
  Let \(m = \frac{1}{2}(\dim(\mathfrak{g}) + \operatorname{rk}(\mathfrak{g}))\) . Every solvable subalgebra of g is of dimension \(\leq m\) ; if it is of dimension m, it is a Borel subalgebra. (Reduce to the algebraically closed case, and use Th. 2.)
\item
  Assume that \(k\) is \(\mathbf{R}, \mathbf{C}\) , or a non-discrete complete ultrametric field. Give the grassmannian \(\mathbf{G}(\mathfrak{g})\) of vector subspaces of \(\mathfrak{g}\) its natural structure of analytic manifold over \(k\) (Differentiable and Analytic Manifolds, Results, 5.2.6). Consider the subsets of \(\mathbf{G}(\mathfrak{g})\) formed by:
\end{enumerate}

\begin{enumerate}
\def\labelenumi{(\roman{enumi})}
\item
  the subalgebras
\item
  the solvable subalgebras
\item
  the nilpotent subalgebras
\item
  the subalgebras consisting of nilpotent elements
\item
  the Borel subalgebras.
\end{enumerate}

Show that these subsets are closed (for (v), use Exerc. 2). Deduce that these subsets are compact when \(k\) is locally compact.

Show by examples that the subsets of \(\mathbf{G}(\mathfrak{g})\) formed by:

\begin{enumerate}
\def\labelenumi{(\roman{enumi})}
\setcounter{enumi}{5}
\item
  the Cartan subalgebras
\item
  the subalgebras reductive in \(\mathfrak{g}\)
\item
  the semi-simple subalgebras
\item
  the decomposable subalgebras
\end{enumerate}

are not necessarily closed, even when \(k = \mathbf{C}\) .

\begin{enumerate}
\def\labelenumi{\arabic{enumi})}
\setcounter{enumi}{3}
\tightlist
\item
  Assume that \(k = \mathbf{C}\) . Let \(G = \operatorname{Int}(\mathfrak{g}) = \operatorname{Aut}_0(\mathfrak{g})\) , and let \(B\) be an integral subgroup of \(G\) whose Lie algebra \(\mathfrak{b}\) is a Borel subalgebra of \(\mathfrak{g}\) . Show that \(B\) is the normalizer of \(\mathfrak{b}\) in \(G\) (use Exerc. 4 of §3 and Exerc. 11 of §5). Deduce by means of Exerc. 3 that \(G / B\) is compact.
\end{enumerate}

¶5) Assume that g is splittable. If h is a splitting Cartan subalgebra of g, denote by \(\mathrm{E}(\mathfrak{h})\) the subgroup of \(\mathrm{Aut}_{e}(\mathfrak{g})\) generated by the \(e^{\operatorname{ad} x}, x \in \mathfrak{g}^{\alpha}(\mathfrak{h}), \alpha \in \mathrm{R}(\mathfrak{g}, \mathfrak{h})\) , cf.~Chap. VII, §3, no. 2.

\begin{enumerate}
\def\labelenumi{\alph{enumi})}
\item
  Let \(\mathfrak{b}\) be a Borel subalgebra of \(\mathfrak{g}\) containing \(\mathfrak{h}\) , and let \(\mathfrak{h}_1\) be a Cartan subalgebra of \(\mathfrak{b}\) . Show that \(\mathfrak{h}_1\) is conjugate to \(\mathfrak{h}\) by an element of \(\mathrm{E}(\mathfrak{h})\) . Deduce that \(\mathrm{E}(\mathfrak{h}) = \mathrm{E}(\mathfrak{h}_1)\) .
\item
  Let \(\mathfrak{h}'\) be a splitting Cartan subalgebra of \(\mathfrak{g}\) . Show that \(\mathrm{E}(\mathfrak{g}) = \mathrm{E}(\mathfrak{h}')\) . (If \(\mathfrak{b}'\) is a Borel subalgebra containing \(\mathfrak{h}'\) , choose a Cartan subalgebra \(\mathfrak{h}_1\) of \(\mathfrak{b} \cap \mathfrak{b}'\) and apply \(a\) ) to show that \(\mathrm{E}(\mathfrak{h}) = \mathrm{E}(\mathfrak{h}_1) = \mathrm{E}(\mathfrak{h}')\) .)
\item
  Let \(x\) be a nilpotent element of \(\mathfrak{g}\) . Show that \(e^{\mathrm{ad}x} \in \mathrm{E}(\mathfrak{h})\) . (Use b) and Cor. 2 of Th. 1 to reduce to the case in which \(x \in [\mathfrak{b}, \mathfrak{b}]\) .) Deduce that \(\mathrm{E}(\mathfrak{h}) = \mathrm{Aut}_e(\mathfrak{g})\) .
\end{enumerate}

\begin{enumerate}
\def\labelenumi{\arabic{enumi})}
\setcounter{enumi}{5}
\tightlist
\item
  \begin{enumerate}
  \def\labelenumii{\alph{enumii})}
  \tightlist
  \item
    Show the equivalence of the properties:
  \end{enumerate}
\end{enumerate}

\begin{enumerate}
\def\labelenumi{(\roman{enumi})}
\item
  g has no nilpotent element \(\neq 0\) .
\item
  g has no parabolic subalgebra ≠ g.
\end{enumerate}

(Use Cor. 2 of Th. 1.)

Such an algebra is called anisotropic.

\begin{enumerate}
\def\labelenumi{\alph{enumi})}
\setcounter{enumi}{1}
\tightlist
\item
  Let \(\mathfrak{p}\) be a minimal parabolic subalgebra of \(\mathfrak{g}\) , \(\mathfrak{r}\) the radical of \(\mathfrak{p}\) , and \(\mathfrak{s} = \mathfrak{p} / \mathfrak{r}\) . Show that \(\mathfrak{s}\) is anisotropic. (Remark that, if \(\mathfrak{q}\) is a parabolic subalgebra of \(\mathfrak{s}\) , the inverse image of \(\mathfrak{q}\) in \(\mathfrak{p}\) is a parabolic subalgebra of \(\mathfrak{g}\) , cf.~§3, Exerc. 5 a).)
\end{enumerate}

¶ 7) a) Show that the following properties of k are equivalent:

\begin{enumerate}
\def\labelenumi{(\roman{enumi})}
\item
  Every anisotropic semi-simple \(k\) -Lie algebra reduces to 0.
\item
  Every semi-simple k-Lie algebra has a Borel subalgebra.
\end{enumerate}

(Use Exerc. 6 to prove that (i) \(\Longrightarrow\) (ii).)

\begin{enumerate}
\def\labelenumi{\alph{enumi})}
\setcounter{enumi}{1}
\tightlist
\item
  Show that (i) and (ii) imply \(^{19}\) :
\end{enumerate}

\begin{enumerate}
\def\labelenumi{(\roman{enumi})}
\setcounter{enumi}{2}
\tightlist
\item
  Every finite dimensional \(k\) -algebra that is a field is commutative. (Or again: the Brauer group of every algebraic extension of \(k\) reduces to 0.)
\end{enumerate}

(Use the Lie algebra of elements of trace zero in such an algebra.)

\begin{enumerate}
\def\labelenumi{\alph{enumi})}
\setcounter{enumi}{2}
\tightlist
\item
  Show that (i) and (ii) are implied by:
\end{enumerate}

\begin{enumerate}
\def\labelenumi{(\roman{enumi})}
\setcounter{enumi}{3}
\tightlist
\item
  For any finite family of homogeneous polynomials \(f_{\alpha} \in k[(X_i)_{i \in \mathrm{I}}]\) of degrees \(\geq 1\) such that \(\sum_{\alpha} \deg f_{\alpha} < \operatorname{Card}(\mathrm{I})\) , there exist elements \(x_i \in k\) , not all zero, such that \(f_{\alpha}((x_i)_{i \in \mathrm{I}}) = 0\) for all \(\alpha\) .
\end{enumerate}

(Use Prop. 5 of §8.)

\exercisesection{\S~11}

\begin{enumerate}
\def\labelenumi{\arabic{enumi})}
\tightlist
\item
  Let \(\mathfrak{g} = \mathfrak{sl}(2, k)\) . Put \(G = \operatorname{Aut}_e(\mathfrak{g})\) ; this group can be identified with \(\mathbf{PSL}_2(k)\) , cf.~Chap. VII, §3, no. 1, Remark 2.
\end{enumerate}

\begin{enumerate}
\def\labelenumi{\alph{enumi})}
\item
  Every nilpotent element of \(\mathfrak{g}\) is G-conjugate to \(\left( \begin{array}{cc}0 & \lambda \\ 0 & 0 \end{array} \right)\) for some \(\lambda \in k\) . Such an element is principal if and only if it is non-zero.
\item
  The elements \(\begin{pmatrix}0&\lambda\\0&0\end{pmatrix},\begin{pmatrix}0&\mu\\0&0\end{pmatrix}\) , where \(\lambda,\mu\in k^{*}\) , are G-conjugate if and only if \(\lambda^{-1}\mu\) is a square in k.
\item
  Every simple element of \(\mathfrak{g}\) is G-conjugate to \(\left( \begin{array}{cc}1 & 0\\ 0 & -1 \end{array} \right)\) .
\end{enumerate}

\begin{enumerate}
\def\labelenumi{\arabic{enumi})}
\setcounter{enumi}{1}
\tightlist
\item
  Let \(A = \begin{pmatrix} 0 & 1 \\ 0 & 0 \end{pmatrix}\) , \(B = \begin{pmatrix} 0 & 0 \\ 1 & 0 \end{pmatrix}\) . Then \(A\) is nilpotent, \(AB\) is not, and
\end{enumerate}

\[
[ A, [ A, B ] ] = \left( \begin{array}{c c} 0 & - 2 \\ 0 & 0 \end{array} \right).
\]

Deduce that Lemma 5 does not extend to fields of characteristic 2.

\begin{enumerate}
\def\labelenumi{\arabic{enumi})}
\setcounter{enumi}{2}
\tightlist
\item
  Let \(\mathfrak{r}\) be the radical of \(\mathfrak{g}\) , and \(\mathfrak{s} = \mathfrak{g} / \mathfrak{r}\) . Show the equivalence of:
\end{enumerate}

\begin{enumerate}
\def\labelenumi{(\roman{enumi})}
\item
  g contains no \(sl_{2}\) -triplet;
\item
  \(\mathfrak{s}\) contains no \(\mathfrak{sl}_2\) -triplet;
\item
  \(\mathfrak{s}\) is anisotropic (§10, Exerc. 6);
\item
  \(\mathfrak{s}\) contains no non-zero diagonalizable element (§3, Exerc. 10).
\end{enumerate}

(Use Prop. 2 and Exerc. 10 a) of §3.)

\begin{enumerate}
\def\labelenumi{\arabic{enumi})}
\setcounter{enumi}{3}
\tightlist
\item
  Let V be a vector space of dimension \(n \geq 2\) , \(\mathfrak{g} = \mathfrak{sl}(V)\) and \(G = \mathbf{PGL}(V)\) , identified with a group of automorphisms of g. An \(sl_{2}\) -triplet in g gives V a faithful \(\mathfrak{sl}(2, k)\) -module structure, and conversely any such structure arises from an \(sl_{2}\) -triplet; an \(sl_{2}\) -triplet is principal if and only if the corresponding \(\mathfrak{sl}(2, k)\) -module is simple; two \(sl_{2}\) -triplets are G-conjugate if and only if the corresponding \(\mathfrak{sl}(2, k)\) -modules are isomorphic. Deduce that the G-conjugacy classes of \(sl_{2}\) -triplets correspond bijectively with the families \((m_{1}, m_{2}, \ldots)\) of integers \(\geq 0\) such that
\end{enumerate}

\[
m _ {1} + 2 m _ {2} + 3 m _ {3} + \dots = n \quad \text { and } \quad m _ {1} <   n.
\]

¶ 5) Assume that g is semi-simple. Let a be a subalgebra of g, reductive in g, of the same rank as g, and containing a principal \(sl_{2}\) -triplet of g. Show that \(a = g\) .

¶6) Assume that g is absolutely simple, and denote its Coxeter number by h. Let x be a nilpotent element of g. Show that \((\operatorname{ad} x)^{2h-1} = 0\) and that \((\operatorname{ad} x)^{2h-2} \neq 0\) if and only if x is principal. (Reduce to the case in which g is split, and x is contained in the subalgebra \(n_{+}\) of Prop. 10. Repeat the proof of Prop. 10.)

¶7) Assume that g is semi-simple. Let x be a nilpotent element of g. Then x is principal if and only if x is contained in a unique Borel subalgebra of g. (Reduce to the case in which k is algebraically closed. Use Prop. 10 and Prop. 10 of §3, no. 3.)

\begin{enumerate}
\def\labelenumi{\arabic{enumi})}
\setcounter{enumi}{7}
\item
  A semi-simple Lie algebra has an \(\mathfrak{sl}_2\) -triplet if and only if it is \(\neq 0\) and has a Borel subalgebra.
\item
  Assume that g is semi-simple. Let N (resp. P) be the set of nilpotent (resp. principal nilpotent) elements of g.
\end{enumerate}

\begin{enumerate}
\def\labelenumi{\alph{enumi})}
\item
  Show that \(\mathrm{P}\) is an open subset of \(\mathbf{N}\) in the Zariski topology (use Exerc. 6).
\item
  Assume that g is splittable. Show that P is dense in N. (Use Prop. 10 and Cor. 2 of Th. 1 of §10.)
\end{enumerate}

\begin{enumerate}
\def\labelenumi{\arabic{enumi})}
\setcounter{enumi}{9}
\tightlist
\item
  Assume that \(\mathfrak{g}\) is splittable semi-simple. Let \((x,h,y)\) be an \(\mathfrak{sl}_2\) -triplet in \(\mathfrak{g}\) .
\end{enumerate}

\begin{enumerate}
\def\labelenumi{\alph{enumi})}
\item
  Show that there exists a splitting Cartan subalgebra \(\mathfrak{h}\) of \(\mathfrak{g}\) containing \(h\) . (Use Exerc. 10 b) of §3.)
\item
  Choose \(\mathfrak{h}\) as in \(a\) ). Show that \(h \in \mathfrak{h}_{\mathbf{Q}}\) and that there exists a basis B of \(\mathrm{R}(\mathfrak{g}, \mathfrak{h})\) such that \(\alpha(h) \in \{0, 1, 2\}\) for all \(\alpha \in \mathrm{B}\) (cf.~Prop. 5). The element \(x\) belongs to the subalgebra of \(\mathfrak{g}\) generated by the \(\mathfrak{g}^{\alpha}\) , \(\alpha \in \mathrm{B}\) .
\item
  Deduce from a) and b), and the Jacobson-Morozov theorem, a new proof of the fact that every nilpotent element of g is contained in a Borel subalgebra (cf.~§10, Cor. 2 of Th. 1).
\end{enumerate}

¶11) Let \((x,h,y)\) be a principal \(\mathfrak{sl}_2\) -triplet in the semi-simple Lie algebra \(\mathfrak{g}\) . Give \(\mathfrak{g}\) the \(\mathfrak{sl}(2,k)\) -module structure defined by this triplet. Show that the module thus defined is isomorphic to \(\bigoplus_{i=1}^{l} V(2k_i - 2)\) , where the \(k_i\) are the characteristic degrees of the algebra of invariant polynomial functions on \(\mathfrak{g}\) . (Reduce to the case in which \(\mathfrak{g}\) is splittable simple. Use Cor. 1 of Th. 1 of §8, no. 3, and Chap. VI, §4, Exerc. 6 c.) \(^{20}\)

¶ 12) Assume that g is semi-simple. Let \(x \in g\) and let s (resp. n) be the semi-simple (resp. nilpotent) component of x. Let \(a_{x}\) (resp. \(a_{s}\) ) be the commutant of x (resp. s) in g.

\begin{enumerate}
\def\labelenumi{\alph{enumi})}
\item
  Show that \(n\) is a nilpotent element of the semi-simple algebra \(\mathcal{D}(\mathfrak{a}_s)\) , and that the commutant of \(n\) in \(\mathfrak{a}_s\) is equal to \(\mathfrak{a}_x\) . Deduce that \(\dim \mathfrak{a}_x < \dim \mathfrak{a}_s\) if \(n \neq 0\) , i.e.~if \(x\) is not semi-simple.
\item
  Show that \(\dim \mathfrak{a}_x = \mathrm{rk}(\mathfrak{g})\) if and only if \(n\) is a principal nilpotent element of \(\mathcal{D}(\mathfrak{a}_s)\) .
\item
  Put \(\mathrm{G} = \mathrm{Aut}_{e}(\mathfrak{g})\) . Show that, for all \(\lambda \in k\) , there exists \(\sigma_{\lambda} \in G\) such that \(\sigma_{\lambda}x = s + \lambda^{2}n\) . (Show that, if \(n \neq 0\) , there exists an \(sl_{2}\) -triplet in \(a_{s}\) of which the first component is n, and deduce a homomorphism \(\varphi : \mathbf{SL}(2, k) \to \mathbf{G}\) ; take \(\sigma_{\lambda}\) to be the image under \(\varphi\) of a suitable diagonal element of \(\mathbf{SL}(2, k)\) .) Deduce that s belongs to the closure of G.x in the Zariski topology.
\end{enumerate}

\(d\) ) Show that, if \(x\) is not semi-simple, \(x\) does not belong to the closure of G.s in the Zariski topology (use the inequality \(\dim \mathfrak{a}_x < \dim \mathfrak{a}_s\) , cf.~\(a\) )).

\begin{enumerate}
\def\labelenumi{\alph{enumi})}
\setcounter{enumi}{4}
\tightlist
\item
  Assume that \(k\) is algebraically closed. Prove the equivalence of the following properties:
\end{enumerate}

\begin{enumerate}
\def\labelenumi{(\roman{enumi})}
\item
  \(x\) is semi-simple;
\item
  G.x is closed in g in the Zariski topology.
\end{enumerate}

(The implication (ii) \(\Longrightarrow\) (i) follows from c). If (i) is satisfied, and if \(x'\) is in the closure of G.x, Exerc. 15 of §8 shows that the semi-simple component \(s'\) of \(x'\) belongs to G.x, so \(x'\) belongs to the closure of G. \(s'\) ; conclude by applying d) to \(x'\) and \(s'\) .)

\(f\) ) Assume that \(k\) is algebraically closed. Let \(\mathrm{F}_x\) be the set of elements \(y\in \mathfrak{g}\) such that \(f(x) = f(y)\) for every invariant polynomial function \(f\) on \(\mathfrak{g}\) ; we have \(y\in \mathrm{F}_x\) if and only if the semi-simple component of \(y\) is G-conjugate to \(s\) (§8, Exerc. 17). Show that \(\mathrm{F}_x\) is the union of a finite number of orbits of G, and that this number is \(\leq 3^{l(x)}\) , where \(l(x)\) is the rank of \(\mathcal{D}(\mathfrak{a}_s)\) . Only one of these orbits is closed: that of \(s\) ; only one is open in \(\mathrm{F}_x\) : that consisting of the elements \(y\in \mathrm{F}_x\) such that \(\dim \mathfrak{a}_y = \operatorname {rk}(\mathfrak{g})\) .

\begin{enumerate}
\def\labelenumi{\arabic{enumi})}
\setcounter{enumi}{12}
\tightlist
\item
  Assume that \(\mathfrak{g}\) is semi-simple.
\end{enumerate}

\begin{enumerate}
\def\labelenumi{\alph{enumi})}
\item
  Let \((x, h, y)\) be a principal \(sl_{2}\) -triplet in g, and let b be the Borel subalgebra containing x (Exerc. 7). Show that b is contained in Im ad x.
\item
  Assume that k is algebraically closed. Show that, for any element z in g, there exist \(x, t \in g\) , with x principal nilpotent, such that \(z = [x, t]\) (apply a) to a Borel subalgebra containing z).
\end{enumerate}

\begin{enumerate}
\def\labelenumi{\arabic{enumi})}
\setcounter{enumi}{13}
\tightlist
\item
  Assume that \(\mathfrak{g}\) is semi-simple. Let \(\mathfrak{p}\) be a parabolic subalgebra of \(\mathfrak{g}\) , and let \(f_{1}\) and \(f_{2}\) be two homomorphisms from \(\mathfrak{g}\) to a finite dimensional Lie algebra. Show that \(f_{1}|\mathfrak{p} = f_{2}|\mathfrak{p}\) implies that \(f_{1} = f_{2}\) . (Reduce to the case in which \(\mathfrak{g}\) is split, then to the case in which \(\mathfrak{g} = \mathfrak{sl}(2,k)\) , and use Lemma 1 of no. 1.)
\end{enumerate}

¶ 15) Assume that g is semi-simple.

\begin{enumerate}
\def\labelenumi{\alph{enumi})}
\item
  Let \(x\) be a nilpotent element of \(\mathfrak{g}\) . Show that \(x\) is contained in \(\operatorname{Im}(\operatorname{ad} x)^2\) (use Prop. 2). Deduce that \((\operatorname{ad} x)^2 = 0\) implies \(x = 0\) .
\item
  Let \((x,h,y)\) be an \(\mathfrak{sl}_2\) -triplet in \(\mathfrak{g}\) , and let \(\mathfrak{s} = kx\oplus kh\oplus ky\) . Prove the equivalence of the following conditions:
\end{enumerate}

\begin{enumerate}
\def\labelenumi{(\roman{enumi})}
\item
  \(\operatorname{Im}(\operatorname{ad} x)^{2}=k.x;\)
\item
  the \(\mathfrak{s}\) -module \(\mathfrak{g} / \mathfrak{s}\) is a sum of simple modules of dimension 1 or 2;
\item
  the only eigenvalues of \(\mathrm{ad}_{\mathfrak{g}}h\) distinct from 0, 1 and -1 are 2 and -2, and their multiplicity is equal to 1.
\end{enumerate}

\begin{enumerate}
\def\labelenumi{\alph{enumi})}
\setcounter{enumi}{2}
\item
  Assume that g is splittable simple. Let h be a splitting Cartan subalgebra of g, B a basis of \(\mathrm{R}(\mathfrak{g},\mathfrak{h})\) , and \(\gamma\) the highest root of \(\mathrm{R}(\mathfrak{g},\mathfrak{h})\) relative to B. Let \((x,h,y)\) be an \(sl_{2}\) -triplet such that \(h\in h_{Q}\) and \(\alpha(h)\geq0\) for all \(\alpha\in B\) (cf.~Prop. 5). Show that conditions (i), (ii), (iii) of b) are satisfied if and only if \(h=H_{\gamma}\) , in which case \(x\in g^{\gamma}\) and \(y\in g^{-\gamma}\) .
\item
  Retain the hypotheses of c), and put \(\mathrm{G} = \mathrm{Aut}_{0}(\mathfrak{g})\) . Show that the \(sl_{2}\) -triplets satisfying (i), (ii) and (iii) are G-conjugate (use Exerc. 10). Show that, if \(x\) is a non-zero nilpotent element of \(\mathfrak{g}\) satisfying (i), the closure of \(G.x\) in the Zariski topology is equal to \(\{0\} \cup G.x\) .
\end{enumerate}

\begin{enumerate}
\def\labelenumi{\arabic{enumi})}
\setcounter{enumi}{15}
\tightlist
\item
  Let \((x,h,y)\) be an \(sl_{2}\) -triplet in g. Show that, if -2 is a square in k, the elements x-y and h are conjugate by an element of \(\operatorname{Aut}_{e}(\mathfrak{g})\) (reduce to the case in which \(\mathfrak{g}=\mathfrak{sl}(2,k)\) ). Deduce that, if g is semi-simple, x-y is semi-simple, and that it is regular if and only if h is regular.
\end{enumerate}

¶ 17) Let \((x, h, y)\) be a principal \(\mathfrak{sl}_2\) -triplet in the semi-simple algebra \(\mathfrak{g}\) . For all \(i \in \mathbf{Z}\) , denote by \(\mathfrak{g}_i\) the eigenspace of ad \(h\) relative to the eigenvalue \(i\) ; we have \(\mathfrak{g} = \bigoplus_{i \in \mathbf{Z}} \mathfrak{g}_i\) , and \(\mathfrak{g}_i = 0\) if \(i\) is odd. The direct sum \(\mathfrak{b}\) of the \(\mathfrak{g}_i\) , \(i \geq 0\) , is a Borel subalgebra of \(\mathfrak{g}\) , and \(\mathfrak{g}_0\) is a Cartan subalgebra.

\begin{enumerate}
\def\labelenumi{\alph{enumi})}
\item
  Show that, for all \(z \in \mathfrak{b}\) , the commutant of \(y + z\) in \(\mathfrak{g}\) is of dimension \(l = \operatorname{rk}(\mathfrak{g})\) .
\item
  Let I be the algebra of invariant polynomial functions on \(\mathfrak{g}\) , and \(\mathrm{P}_1, \ldots, \mathrm{P}_l\) homogeneous elements of I generating I; put \(\deg(\mathrm{P}_i) = k_i = m_i + 1\) . Show that the commutant \(\mathfrak{c}\) of \(x\) in \(\mathfrak{g}\) has a basis \(x_1, \ldots, x_l\) with \(x_i \in \mathfrak{g}_{2m_i}\) (cf.~Exerc. 11).
\item
  Let \(i \in \{1, l\}\) , and let \(J_i\) (resp. \(K_i\) ) be the set of \(j \in (1, l)\) such that \(m_j = m_i\) (resp. \(m_j < m_i\) ). Let \(f_i \in k[X_1, \ldots, X_l]\) be the polynomial such that
\end{enumerate}

\[
f _ {i} \left(a _ {1}, \dots , a _ {l}\right) = \mathrm{P} _ {i} \left(y + \sum_ {j = 1} ^ {j = l} a _ {j} x _ {j}\right) \quad \text { for } \left(a _ {j}\right) \in k ^ {l}.
\]

Show that \(f_{i}\) is the sum of a linear form \(\mathrm{L}_i\) in the \(\mathrm{X}_j\) , \(j \in \mathrm{J}_i\) , and a polynomial in the \(\mathrm{X}_j\) , \(j \in \mathrm{K}_i\) . (If \(t \in k^*\) , the automorphism of \(\mathfrak{g}\) equal to \(t^i\) on \(\mathfrak{g}_i\) belongs to \(\operatorname{Aut}_0(\mathfrak{g})\) and takes \(y\) to \(t^{-2}y\) and \(x_j\) to \(t^{2m_j}x_j\) . Use the invariance of \(\mathrm{P}_i\) under this automorphism.)

\(d)\) Let \(\mathrm{P}\) be the map from \(\mathfrak{g}\) to \(k^l\) defined by the \(\mathrm{P}_i\) . If \(z \in \mathfrak{g}\) , denote by \(\mathrm{D}_z\mathrm{P} : \mathfrak{g} \to k^l\) the tangent linear map of \(\mathrm{P}\) at \(z\) (Chap. VII, App. I, no. 2).

Show that \(\mathfrak{c}\cap\operatorname{Im}\operatorname{ad}(y-x)=0\) (decompose g into a direct sum of simple submodules relative to the subalgebra generated by the given \(sl_{2}\) -triplet). Deduce that the restriction of \(D_{y-x}P\) to c is an isomorphism from c to \(k^{l}\) (use the preceding exercise and Exerc. 18 a) of §8). Show, by using this result, that the determinant of the linear forms \(L_{i}\) defined in c) is \(\neq0\) , and deduce the following results:

\(d_{1})\) the polynomials \(f_{1},\ldots ,f_{l}\) are algebraically independent and generate \(k[\mathrm{X}_1,\dots ,\mathrm{X}_l]\) ;

\(d_{2})\) the map \(z \mapsto \mathrm{P}(y + z)\) from \(\mathfrak{c}\) to \(k^{l}\) is bijective polynomial, and the inverse map is polynomial;

\(d_{3})\) for all \(z\in y + \mathfrak{c}\) , the linear map \(\mathrm{D}_z\mathrm{P}|\mathfrak{c}\) is of rank \(l\) .

In particular, the map \(\mathrm{P}:\mathfrak{g}\to k^l\) is surjective.

\begin{enumerate}
\def\labelenumi{\alph{enumi})}
\setcounter{enumi}{4}
\tightlist
\item
  If \(k\) is algebraically closed, every element of \(\mathfrak{g}\) whose commutant is of dimension \(l\) is conjugate under \(\mathrm{Aut}_0(\mathfrak{g})\) to a unique element of \(y + \mathfrak{c}\) .
\end{enumerate}

\(f\) ) Give an example of a simple Lie algebra with no principal \(\mathfrak{sl}_2\) -triplet, for which the map \(\mathrm{P}\) is not surjective (take \(k = \mathbf{R}\) and \(l = 1\) ).

\exercisesection{\S~13}

\begin{enumerate}
\def\labelenumi{\arabic{enumi})}
\tightlist
\item
  The dimensions \(\leq 80\) of splittable simple Lie algebras are:
\end{enumerate}

3 (A \(_{1}\) = B \(_{1}\) = C \(_{1}\) ), 8 (A \(_{2}\) ), 10 (B \(_{2}\) = C \(_{2}\) ), 14 (G \(_{2}\) ), 15 (A \(_{3}\) = D \(_{3}\) ), 21 (B \(_{3}\) and C \(_{3}\) ), 24 (A \(_{4}\) ), 28 (D \(_{4}\) ), 35 (A \(_{5}\) ), 36 (B \(_{4}\) and C \(_{4}\) ), 45 (D \(_{5}\) ), 48 (A \(_{6}\) ), 52 (F \(_{4}\) ), 55 (B \(_{5}\) and C \(_{5}\) ), 63 (A \(_{7}\) ), 66 (D \(_{6}\) ), 78 (B \(_{6}\) , C \(_{6}\) and E \(_{6}\) ), 80 (A \(_{8}\) ).

\begin{enumerate}
\def\labelenumi{\arabic{enumi})}
\setcounter{enumi}{1}
\tightlist
\item
  Let \((\mathfrak{g},\mathfrak{h})\) be a split simple Lie algebra, and B a basis of \(R(\mathfrak{g},\mathfrak{h})\) . The simple \(\mathfrak{g}\) -modules \(E(\lambda)\) , \(\lambda \in P_{++} - \{0\}\) , of minimum dimension are those for which \(\lambda\) is one of the following weights:
\end{enumerate}

\(\varpi_{1}\) (A \(_{1}\) ); \(\varpi_{1}\) and \(\varpi_{l}\) (A \(_{l}\) , \(l \geq 2\) ); \(\varpi_{1}\) (B \(_{l}\) and C \(_{l}\) , \(l \geq 2\) ); \(\varpi_{1}\) , \(\varpi_{3}\) and \(\varpi_{4}\) (D \(_{4}\) ); \(\varpi_{1}\) (D \(_{l}\) , \(l \geq 5\) ); \(\varpi_{1}\) and \(\varpi_{6}\) (E \(_{6}\) ); \(\varpi_{7}\) (E \(_{7}\) ); \(\varpi_{8}\) (E \(_{8}\) ); \(\varpi_{4}\) (F \(_{4}\) ); \(\varpi_{1}\) (G \(_{2}\) ). Any two such modules can be transformed into each other by an automorphism of g.

Type \(E_{8}\) is the only one for which the adjoint representation is of minimum dimension.

\begin{enumerate}
\def\labelenumi{\arabic{enumi})}
\setcounter{enumi}{2}
\tightlist
\item
  \begin{enumerate}
  \def\labelenumii{\alph{enumii})}
  \tightlist
  \item
    Define an isomorphism from \(\mathfrak{sl}(4,k)\) to the orthogonal algebra \(\mathfrak{o}_S(6,k)\) . (Use the fact that the representation \(\bigwedge^2\sigma\) of no. 1.V is orthogonal and of dimension 6.) The two types of irreducible representation of \(\mathfrak{sl}(4,k)\) of degree 4 correspond to the two semi-spinor representations of \(\mathfrak{o}_S(6,k)\) .
  \end{enumerate}
\end{enumerate}

\begin{enumerate}
\def\labelenumi{\alph{enumi})}
\setcounter{enumi}{1}
\item
  Define an isomorphism from \(\mathfrak{sp}(4,k)\) to the orthogonal algebra \(\mathfrak{o}_S(5,k)\) . (Use the fact that the representation \(\sigma_2\) of no. 3.V is orthogonal and of dimension 5.) The irreducible representation of \(\mathfrak{sp}(4,k)\) of degree 4 corresponds to the spinor representation of \(\mathfrak{o}_S(5,k)\) .
\item
  Define an isomorphism \(\mathfrak{sl}(2,k)\times\mathfrak{sl}(2,k)\to\mathfrak{o}_{S}(4,k)\) by using the tensor product of the identity representations of the two \(\mathfrak{sl}(2,k)\) factors. Recover this result by means of Chap. I, §6, Exerc. 26.
\end{enumerate}

\begin{enumerate}
\def\labelenumi{\arabic{enumi})}
\setcounter{enumi}{3}
\tightlist
\item
  Let \(S\) be the square matrix of order \(n\)
\end{enumerate}

\[
(\delta_ {i, n + 1 - i}) = \left( \begin{array}{c c c c c} 0 & 0 & \dots & 0 & 1 \\ 0 & 0 & \dots & 1 & 0 \\ \vdots & \vdots & \ddots & \vdots & \vdots \\ 0 & 1 & \dots & 0 & 0 \\ 1 & 0 & \dots & 0 & 0 \end{array} \right).
\]

The elements of \(\mathfrak{o}_{S}(n,k)\) are the matrices \((a_{ij})\) that are anti-symmetric with respect to the second diagonal:

\[
a _ {i, j} = - a _ {n + 1 - j, n + 1 - i} \quad \text {   for   every   pair   } (i, j).
\]

The algebra \(\mathfrak{o}_S(n,k)\) is splittable simple of type \(\mathrm{D}_{n/2}\) if \(n\) is even and \(\geq 6\) , and of type \(\mathrm{B}_{(n-1)/2}\) if \(n\) is odd and \(\geq 5\) . The diagonal (resp. upper triangular) elements of \(\mathfrak{o}_S(n,k)\) form a splitting Cartan (resp. Borel) subalgebra.

\begin{enumerate}
\def\labelenumi{\arabic{enumi})}
\setcounter{enumi}{4}
\item
  The notations are those of no. 1.IV, type \(\mathrm{A}_l\) . Show that, if \(n\) is \(\geq 0\) , \(\mathbf{S}^n\sigma\) is an irreducible representation of \(\mathfrak{sl}(l + 1,k)\) of highest weight \(n\varpi_1\) . (Realise \(\mathbf{S}^n\sigma\) on the space of homogeneous polynomials of degree \(n\) in \(\mathrm{X}_0,\dots ,\mathrm{X}_l\) , and observe that, up to homothety, the only polynomial \(f\) such that \(\partial f / \partial \mathrm{X}_i = 0\) for \(i\geq 1\) is \(\mathrm{X}_0^n\) .) Show that all the weights of this representation are of multiplicity 1.
\item
  The notations are those of no. 2.IV, type \(\mathrm{B}_l\) . Show that, if \(1 \leq r \leq l - 1\) , the dimension of \(\operatorname{E}(\varpi_r)\) is \(\binom{2l + 1}{r}\) , and deduce another proof of the fact that \(\bigwedge^r \sigma\) is a fundamental representation of highest weight \(\varpi_r\) .
\item
  The notations are those of no. 2.VII, type \(\mathrm{B}_l\) . Show that \(\mathbf{O}_0^+ (\Psi)\) is the commutator subgroup of \(\mathbf{SO}(\Psi)\) . (Remark that \(\mathbf{O}(\Psi)\) is equal to \(\{\pm 1\} \times \mathbf{SO}(\Psi)\) , hence has the same commutator subgroup as \(\mathbf{SO}(\Psi)\) , and apply Algebra, Chap. IX, §9, Exerc. 11 b).) Deduce that \(\operatorname{Aut}_e(\mathfrak{g}) = \mathbf{O}_0^+ (\Psi)\) .
\item
  The notations are those of no. 3.IV, type \(C_{l}\) ( \(l \geq 1\) ). Show that \(S^{2}\sigma\) is equivalent to the adjoint representation of g.
\item
  The notations are those of no. 3.VII, type \(C_l\) ( \(l \geq 1\) ). In particular, we identify \(\mathrm{Aut}_e(\mathfrak{g})\) with a subgroup of \(\mathbf{Sp}(\varPsi)/\{\pm 1\}\) . Show that the image in \(\mathbf{Sp}(\varPsi)/\{\pm 1\}\) of a symplectic transvection (Algebra, Chap. IX, §4, Exerc. 6) belongs to \(\mathrm{Aut}_e(\mathfrak{g})\) . Deduce that \(\mathrm{Aut}_e(\mathfrak{g}) = \mathbf{Sp}(\varPsi)/\{\pm 1\}\) (Algebra, Chap. IX, §5, Exerc. 11), and that \(\mathrm{Aut}(\mathfrak{g})/\mathrm{Aut}_e(\mathfrak{g})\) can be identified with \(k^*/k^{*2}\) .
\end{enumerate}

¶ 10) The notations are those of no. 4.IV, type D \(_{l}\) (l ≥ 2).

\begin{enumerate}
\def\labelenumi{\alph{enumi})}
\tightlist
\item
  Let \(x\) and \(y\) be the elements of \(\bigwedge^{l} \mathrm{V}\) defined by
\end{enumerate}

\[
x = e _ {1} \wedge \dots \wedge e _ {l - 1} \wedge e _ {l} \quad \text { and } \quad y = e _ {1} \wedge \dots \wedge e _ {l - 1} \wedge e _ {- l}.
\]

The element x is primitive of weight \(2\varpi_{l}\) and y is primitive of weight \(2\varpi_{l-1}\) . The submodule X (resp. Y) of \(\bigwedge^{l}V\) generated by x (resp. y) is isomorphic to \(\mathrm{E}(2\varpi_{l})\) (resp. \(\mathrm{E}(2\varpi_{l-1})\) ). Show, by calculating dimensions, that \(\bigwedge^{l}V = X \oplus Y\) ; in particular, \(\bigwedge^{l}V\) is the sum of two non-isomorphic simple modules.

\begin{enumerate}
\def\labelenumi{\alph{enumi})}
\setcounter{enumi}{1}
\tightlist
\item
  Let \(e = e_{1} \wedge \cdots \wedge e_{l} \wedge e_{-1} \wedge \cdots \wedge e_{-l} \in \bigwedge^{2l}V\) , and \(\varPsi_{l}\) the extension of \(\varPsi\) to \(\bigwedge^{l}V\) . Let \(z \in \bigwedge^{l}V\) . Prove the equivalences:
\end{enumerate}

\[
z \in X \Longleftrightarrow z \wedge t = \Psi_ {l} (z, t) e \quad \text {   for   all   } t \in \bigwedge^ {l} V
\]

\[
z \in \mathrm{Y} \Longleftrightarrow z \wedge t = - \Psi_ {l} (z, t) e \quad \text {   for   all   } t \in \bigwedge^ {l} \mathrm{V}.
\]

(If \(X'\) and \(Y'\) denote the subspaces defined by the right-hand sides, prove first that \(X'\) and \(Y'\) are stable under \(\mathfrak{g}\) and contain \(x\) and \(y\) , respectively.)

\begin{enumerate}
\def\labelenumi{\alph{enumi})}
\setcounter{enumi}{2}
\item
  Assume that z is pure (Algebra, Chap. III, §11, no. 13), and denote by \(M_{z}\) the l-dimensional subspace of V associated to it. Show that \(M_{z}\) is totally isotropic if and only if z belongs to X or to Y. (Use the fact that, when \(M_{z}\) is totally isotropic, there exists an orthogonal transformation that transforms it into \(M_{x}\) ; when \(M_{z}\) is not totally isotropic, construct an l-vector t such that \(z \wedge t = 0, \Psi_{l}(z, t) = 1\) , and apply b) above.) When \(z \in X\) (resp. \(z \in Y\) ), the dimension of \(\mathrm{M}_{x}/(\mathrm{M}_{x} \cap \mathrm{M}_{z})\) is an even (resp. odd) integer, cf.~Algebra, Chap. IX, §6, Exerc. 18 d).
\item
  Let s be a direct (resp. inverse) similarity of V. Show that \(\bigwedge^{l}s\) leaves X and Y stable (resp. interchanges X and Y).
\end{enumerate}

\begin{enumerate}
\def\labelenumi{\arabic{enumi})}
\setcounter{enumi}{10}
\item
  The notations are those of no. 4.VII, type \(\mathrm{D}_l\) ( \(l \geq 3\) ). Show that \(\mathbf{O}_0^+(\varPsi)\) is the commutator subgroup of \(\mathbf{SO}(\varPsi)\) . (Apply Algebra, Chap. IX, §6, Exerc. 17 b) and Algebra, Chap. IX, §9, Exerc. 11 b). Deduce that \(\mathrm{Aut}_e(\mathfrak{g})\) is equal to the image of \(\mathbf{O}_0^+(\varPsi)\) in \(\mathbf{SO}(\varPsi)/\{\pm 1\}\) . Moreover, \(-1\) belongs to \(\mathbf{O}_0^+(\varPsi)\) if and only if \(l\) is even or \(-1\) is a square in \(k\) (Algebra, Chap. IX, §9, Exerc. 11 c)).
\item
  The Killing form of \(\mathfrak{sl}(n,k)\) is \((X,Y)\mapsto 2n\operatorname {Tr}(XY)\) . That of \(\mathfrak{sp}(n,k)\) , \(n\) even, is \((X,Y)\mapsto (n + 2)\operatorname {Tr}(XY)\) . That of \(\mathfrak{o}_S(n,k)\) , where \(S\) is non-degenerate symmetric of rank \(n\) , is \((X,Y)\mapsto (n - 2)\operatorname {Tr}(XY)\) .
\item
  The algebra of invariant polynomial functions of \(\mathfrak{g}\) is generated:
\end{enumerate}

\(a\) ) in case \(\mathrm{A}_l\) , by the functions \(X\mapsto \mathrm{Tr}(X^i)\) , \(2\leq i\leq l + 1\) ;

\begin{enumerate}
\def\labelenumi{\alph{enumi})}
\setcounter{enumi}{1}
\item
  in case \(B_{l}\) , by the functions \(X \mapsto \operatorname{Tr}(X^{2i})\) , \(1 \leq i \leq l\) ;
\item
  in case \(C_{l}\) , by the functions \(X \mapsto \operatorname{Tr}(X^{2i})\) , \(1 \leq i \leq l\) ;
\end{enumerate}

\(d\) ) in case \(\mathrm{D}_l\) , by the functions \(X \mapsto \operatorname{Tr}(X^{2i})\) , \(1 \leq i \leq l - 1\) , and by one of the two polynomial functions \(\tilde{f}\) such that \(\tilde{f}(X)^2 = (-1)^l \det(X)\) .

\begin{enumerate}
\def\labelenumi{\arabic{enumi})}
\setcounter{enumi}{13}
\tightlist
\item
  \begin{enumerate}
  \def\labelenumii{\alph{enumii})}
  \tightlist
  \item
    Let G be the group associated to \(\mathfrak{g} = \mathfrak{sl}(n, k)\) by the procedure of §7, Exerc. 26. The natural g-module structure on \(k^{n}\) gives rise to a homomorphism \(\varphi : G \to \mathbf{GL}(n, k)\) . Use loc. cit. h) to prove that \(\varphi\) is injective, and loc. cit. f) to prove that
  \end{enumerate}
\end{enumerate}

\[
\operatorname{Im} (\varphi) = \mathbf {S L} (n, k).
\]

\begin{enumerate}
\def\labelenumi{\alph{enumi})}
\setcounter{enumi}{1}
\item
  Let E be a finite dimensional \(\mathfrak{sl}(n,k)\) -module, and \(\rho\) the corresponding representation of \(\mathfrak{sl}(n,k)\) . Show that there exists a unique representation \(\pi:\mathbf{SL}(n,k)\to\mathbf{GL}(\mathrm{E})\) such that \(\pi(e^{x})=e^{\rho(x)}\) for every nilpotent element x of \(\mathfrak{sl}(n,k)\) (use a)). We say that \(\rho\) and \(\pi\) are compatible. Generalize the results proved for n=2 in §1, no. 4.
\item
  Assume that \(k\) is \(\mathbf{R}, \mathbf{C}\) or a complete field for a discrete valuation with residue field of characteristic \(\neq 0\) . Show that \(\rho\) and \(\pi\) are compatible if and only if \(\pi\) is a homomorphism of Lie groups such that \(L(\pi) = \rho\) (same method as in Exerc. 18 b) of §1).
\end{enumerate}

\(d\) ) Prove analogous results for \(\mathfrak{sp}(2n,k)\) and \(\mathbf{Sp}(2n,k)\) .

¶15) Let V be a vector space of finite dimension \(\geq 2\) , g a Lie subalgebra of End(V) and \(\theta\) an element of g. We make the following assumptions:

\begin{enumerate}
\def\labelenumi{(\roman{enumi})}
\item
  V is a semi-simple g-module;
\item
  \(\theta\) is of rank 1 (i.e.~\(\dim \operatorname{Im}(\theta) = 1\) );
\item
  the line \(\operatorname{Im}(\theta)\) generates the U(g)-module V.
\end{enumerate}

\begin{enumerate}
\def\labelenumi{\alph{enumi})}
\item
  Show that these assumptions are satisfied (for a suitable choice of \(\theta\) ) when \(\mathfrak{g} = \mathfrak{sl}(V)\) , when \(\mathfrak{g} = \mathfrak{gl}(V)\) , or when there exists a non-degenerate alternating bilinear form \(\varPsi\) on V such that \(\mathfrak{g} = \mathfrak{sp}(\varPsi)\) or \(\mathfrak{g} = k.1 \oplus \mathfrak{sp}(\varPsi)\) . In each of these cases, \(\theta\) can be taken to be a nilpotent element; in the second case (and only in that case) \(\theta\) can be taken to be a semi-simple element.
\item
  We shall prove that the four cases above are the only ones possible. We are reduced immediately to the case in which k is algebraically closed. Show that V is then a simple g-module, and that \(g = c \oplus s\) , where s is semi-simple, and c = 0 or c = k.1; show that V is not isomorphic to a tensor product of g-modules of dimension \(\geq 2\) ; deduce that s is simple.
\item
  Choose a Cartan subalgebra \(\mathfrak{h}\) of \(\mathfrak{s}\) , and a basis B of \(\mathrm{R}(\mathfrak{s},\mathfrak{h})\) . Let \(\lambda\) be the highest weight (with respect to B) of the \(\mathfrak{s}\) -module V, and let \(e\) be a non-zero element of V of weight \(\lambda\) . The highest weight of the dual module \(\mathrm{V}^*\) is \(\lambda^{*} = -w_{0}\lambda\) (§7, no. 5); let \(e^*\) be a non-zero element of \(\mathrm{V}^*\) of weight \(\lambda^*\) . Identify \(\mathrm{V} \otimes \mathrm{V}^*\) with \(\operatorname{End}(\mathrm{V})\) as usual. Show that there exist \(x \in \mathrm{V}, y \in \mathrm{V}^*\) such that \(x \otimes y \in \mathfrak{g}\) and \(\langle x, e^* \rangle \neq 0, \langle e, y \rangle \neq 0\) (take the conjugate of \(\theta\) by \(e^n\) , where \(n\) is a suitable nilpotent element of \(\mathfrak{g}\) ). Use the fact that \(\mathfrak{g}\) is an \(\mathfrak{h}\) -submodule of \(\mathrm{V} \otimes \mathrm{V}^*\) to conclude that \(\mathfrak{g}\) contains \(e \otimes e^*\) . Deduce that \(\lambda + \lambda^{*} = \tilde{\alpha}\) , where \(\tilde{\alpha}\) is the highest root of \(\mathfrak{s}\) .
\end{enumerate}

\(d\) ) Show that \(\mathfrak{s}\) is not of type \(\mathrm{B}_l\) ( \(l\geq 2\) ), \(\mathrm{D}_l\) ( \(l\geq 4\) ), \(\mathrm{E}_6,\mathrm{E}_7,\mathrm{E}_8,\mathrm{F}_4,\mathrm{G}_2\) (by Chap. VI, Tables, \(\tilde{\alpha}\) would be a fundamental weight, and hence could not be of the form \(\lambda +\lambda^{*}\) above). Deduce that \(\mathfrak{s}\) is either of type \(\mathrm{A}_l\) or of type \(\mathrm{C}_l\) , and that in the first case \(\lambda = \varpi_1\) or \(\varpi_1 = \varpi_1^*\) , and in the second case \(\tilde{\alpha} = 2\varpi_1 = \varpi_1 + \varpi_1^*\) ; since \(\mathfrak{c} = 0\) or \(k.1\) , this indeed gives the four possibilities in \(a)^{21}\) .

¶16) Let g be an absolutely simple Lie algebra of type \(A_{l}\) ( \(l \geq 2\) ), \(\bar{k}\) an algebraic closure of k, and \(\pi : \operatorname{Gal}(\bar{k}/k) \to \operatorname{Aut}(\bar{\mathbf{R}},\bar{\mathbf{B}})\) the homomorphism defined in Exerc. 8 of §5.

\begin{enumerate}
\def\labelenumi{\alph{enumi})}
\tightlist
\item
  Assume that \(\pi\) is trivial. Show that there exist exactly two two-sided ideals \(\mathfrak{m}\) and \(\mathfrak{m}'\) of \(\mathrm{U}(\mathfrak{g})\) such that \(\mathrm{D} = \mathrm{U}(\mathfrak{g}) / \mathfrak{m}\) and \(\mathrm{D}' = \mathrm{U}(\mathfrak{g}) / \mathfrak{m}'\) are central simple algebras of dimension \((l + 1)^2\) (use Exerc. 8 of §7). The principal anti-automorphism of \(\mathrm{U}(\mathfrak{g})\) interchanges \(\mathfrak{m}\) and \(\mathfrak{m}'\) ; in particular, \(\mathrm{D}'\) is isomorphic to the opposite of D. The composite \(\mathfrak{g} \to \mathrm{U}(\mathfrak{g}) \to \mathrm{D}\) identifies \(\mathfrak{g}\) with the Lie subalgebra \(\mathfrak{sl}_{\mathrm{D}}\) of D consisting of the elements of trace zero. Moreover, \(\mathfrak{g}\) is splittable (and hence isomorphic to \(\mathfrak{sl}(l + 1,k)\) ) if and only if D is isomorphic to \(\mathbf{M}_{l + 1}(k)\) .
\end{enumerate}

Conversely, if \(\Delta\) is a central simple algebra of dimension \((l + 1)^2\) , the Lie algebra \(\mathfrak{sl}_{\Delta}\) is absolutely simple of type \(A_l\) , and the corresponding homomorphism \(\pi\) is trivial. Two such algebras \(\mathfrak{sl}_{\Delta}\) and \(\mathfrak{sl}_{\Delta'}\) are isomorphic if and only if \(\Delta\) and \(\Delta'\) are isomorphic or anti-isomorphic.

\begin{enumerate}
\def\labelenumi{\alph{enumi})}
\setcounter{enumi}{1}
\tightlist
\item
  Assume that \(\pi\) is non-trivial. Since \(\operatorname{Aut}(\bar{\mathbf{R}},\bar{\mathbf{B}})\) has two elements, the kernel of \(\pi\) is an open subgroup of \(\operatorname{Gal}(\bar{k}/k)\) of index 2, which corresponds by Galois theory to a quadratic extension \(k_{1}\) of k. Denote the non-trivial involution of \(k_{1}\) by \(x\mapsto\bar{x}\) .
\end{enumerate}

Show that there exists a unique two-sided ideal \(\mathfrak{m}\) of \(\mathrm{U}(\mathfrak{g})\) such that \(\mathrm{D} = \mathrm{U}(\mathfrak{g}) / \mathfrak{m}\) is a simple algebra of dimension \(2(l + 1)^2\) whose centre is a quadratic extension of \(k\) (same method). The centre of D can be identified with \(k_{1}\) . The ideal \(\mathfrak{m}\) is stable under the principal anti-automorphism of \(\mathrm{U}(\mathfrak{g})\) ; that defines by passage to the quotient an involutive anti-automorphism \(\sigma\) of D such that \(\sigma(x) = \bar{x}\) for all \(x \in k_{1}\) . The composite map \(\mathfrak{g} \to \mathrm{U}(\mathfrak{g}) \to \mathrm{D}\) identifies \(\mathfrak{g}\) with the Lie subalgebra \(\mathfrak{s}\mathfrak{u}_{\mathrm{D},\sigma}\) of D consisting of the elements \(x\) such that \(\sigma(x) = -x\) and \(\mathrm{Tr}_{\mathrm{D}/k_1}(x) = 0\) . Conversely, if \(\Delta\) is a central simple \(k_{1}\) -algebra of dimension \((l + 1)^2\) , equipped with an involutive anti-automorphism \(\sigma\) whose restriction to \(k_{1}\) is \(x \mapsto \bar{x}\) , the Lie algebra \(\mathfrak{s}\mathfrak{u}_{\Delta,\sigma}\) is absolutely simple of type \(\mathrm{A}_l\) , and the corresponding homomorphism \(\pi\) is that associated to \(k_{1}\) . Two such algebras \(\mathfrak{s}\mathfrak{u}_{\Delta,\sigma}\) and \(\mathfrak{s}\mathfrak{u}_{\Delta',\sigma'}\) are isomorphic if and only if there exists a \(k\) -isomorphism \(f: \Delta \to \Delta'\) such that \(\sigma' \circ f = f \circ \sigma\) .

Show that, when \(\mathrm{D} = \mathbf{M}_{l+1}(k_{1})\) , there exists an invertible hermitian matrix H of degree \(l + 1\) , unique up to multiplication by an element of \(k^{*}\) , such that

\[
\sigma (x) = H. ^ {t} \bar {x}. H ^ {- 1}
\]

for all \(x \in \mathbf{M}_{l+1}(k_1)\) ; the algebra g can then be identified with the algebra \(\mathfrak{su}(l + 1, H)\) consisting of the matrices x such that \(x.H + H.^t\bar{x} = 0\) and \(\operatorname{Tr}(x) = 0\) .

¶ 17) Let g be an absolutely simple Lie algebra of type \(B_{l}\) (resp. \(C_{l}, D_{l}\) ), with \(l \geq 2\) (resp. \(l \geq 3, l \geq 4\) ). Assume further that, when g is of type \(D_{4}\) , the image of the homomorphism \(\pi : \operatorname{Gal}(\bar{k}/k) \to \operatorname{Aut}(\bar{\mathrm{R}}, \bar{\mathrm{B}})\) defined in Exerc. 8 of §5 is of order \(\leq 2\) . Show that there exists a central simple algebra D of dimension \((2l + 1)^{2}\) (resp. \(4l^{2}, 4l^{2}\) ) and an involutive anti-automorphism \(\sigma\) of D such that g is isomorphic to the Lie subalgebra of D consisting of the elements x such that \(\sigma(x) = -x\) and \(\operatorname{Tr}_{\mathrm{D}}(x) = 0\) (same method as in Exerc. 16 a)). \(^{22}\)

\begin{enumerate}
\def\labelenumi{\arabic{enumi})}
\setcounter{enumi}{17}
\tightlist
\item
  Let V be a finite dimensional vector space, Q a non-degenerate quadratic form on V, and \(\Psi\) the symmetric bilinear form associated to Q. Denote the Lie algebra \(\mathfrak{o}(\varPsi)\) by g and the extension of \(\varPsi\) to \(\bigwedge^{2}(V)\) by \(\varPsi_{2}\) .
\end{enumerate}

\begin{enumerate}
\def\labelenumi{\alph{enumi})}
\tightlist
\item
  Show that there exists an isomorphism of vector spaces \(\theta : \bigwedge^{2}(\mathrm{V}) \to \mathfrak{g}\) characterized by the following equivalent properties:
\end{enumerate}

\begin{enumerate}
\def\labelenumi{(\roman{enumi})}
\item
  For \(a, b\) and \(x\) in \(\mathrm{V}\) , we have \(\theta(a \wedge b). x = a. \Psi(x, b) - b. \Psi(x, a)\) .
\item
  For \(x, y\) in V and \(u\) in \(\bigwedge^2 (\mathrm{V})\) we have \(\varPsi_2(x\wedge y,u) = \varPsi (x,\theta (u).y)\) .
\end{enumerate}

Let \(\sigma\) be the identity representation of \(\mathfrak{g}\) on \(V\) ; then \(\theta\) is an isomorphism of \(\bigwedge^{2}(V)\) with the adjoint representation of \(\mathfrak{g}\) .

\begin{enumerate}
\def\labelenumi{\alph{enumi})}
\setcounter{enumi}{1}
\item
  Define the linear map \(f: \mathfrak{g} \to \mathrm{C}^{+}(\mathrm{Q})\) as in Lemma 1 of no. 2. Show that \(f\theta(a \wedge b) = \frac{1}{2}(ab - ba)\) for a, b in V, and deduce a new proof of assertions (iii), (iv) and (v) of Lemma 1 of no. 2.
\item
  The notations \(l, \mathrm{F}, \mathrm{F}', e_0, \mathrm{N}, \lambda\) and \(\rho\) are those of no. 2. Choose \(e \neq 0\) in \(\bigwedge (\mathrm{F}')\) and define the bilinear form \(\varPhi\) on \(\mathrm{N}\) by
\end{enumerate}

\[
x \wedge y = (- 1) ^ {p (p + 1) / 2} \Phi (x, y). e \quad (x \in \bigwedge^ {p} (\mathrm{F} ^ {\prime}), y \in \mathrm{N}).
\]

Show that \(\Phi(\lambda(a).x,y)+\Phi(x,\lambda(a).y)=0\) for x,y in N and a in \(F\oplus F'\) . Deduce that \(\Phi\) is invariant under the representation \(\rho\) of g (remark that the Lie algebra \(\rho(\mathfrak{g})\) is generated by \(\lambda(F\oplus F')\) by a) and b) above).

\begin{enumerate}
\def\labelenumi{\arabic{enumi})}
\setcounter{enumi}{18}
\tightlist
\item
  Let \(\mathrm{V}\) be a finite dimensional vector space and \(\mathrm{E} = \mathrm{V} \oplus \mathrm{V}^*\) . Define a non-degenerate bilinear form \(\varPhi\) on \(\mathrm{E}\) by
\end{enumerate}

\[
\Phi ((x, x ^ {*}), (y, y ^ {*})) = \langle x, y ^ {*} \rangle + \langle y, x ^ {*} \rangle .
\]

Put \(N = \bigwedge(V)\) and \(Q(x) = \frac{1}{2}\Phi(x, x)\) for \(x \in E\) . Denote the spinor representation by \(\lambda : C(Q) \to \text{End}(N)\) as in no. 2.IV, define \(f : \mathfrak{o}(\Phi) \to C^{+}(Q)\) as in Lemma 1 of no. 2, and denote by \(\rho\) the linear representation \(\lambda \circ f\) of the algebra \(\mathfrak{o}(\Phi)\) on N.

\begin{enumerate}
\def\labelenumi{\alph{enumi})}
\item
  Associate to any endomorphism \(u\) of \(V\) the endomorphism \(\tilde{u}\) of \(E\) by the formula \(\tilde{u}(x, x^{*}) = (u(x), -^{t}u(x^{*}))\) . Show that \(u \mapsto \tilde{u}\) is a homomorphism of Lie algebras from \(\mathfrak{gl}(V)\) to \(\mathfrak{o}(\varPhi)\) ; moreover, for \(u\) in \(\mathfrak{gl}(V)\) , \(\rho(\tilde{u})\) is the unique derivation of the algebra \(\bigwedge(V) = N\) that coincides with \(u\) on \(V\) .
\item
  Let \(\varPsi\) be a non-degenerate alternating bilinear form on V and \(\gamma: V \to V^{*}\) the isomorphism defined by \(\varPsi(x, y) = \langle x, \gamma(y) \rangle\) for x, y in V. Show that the endomorphisms \(\tilde{X}_{+}\) and \(\tilde{X}_{-}\) of E defined by
\end{enumerate}

\[
\tilde {X} _ {+} (x, x ^ {*}) = (\gamma^ {- 1} (x ^ {*}), 0), \quad \tilde {X} _ {-} (x, x ^ {*}) = (0, - \gamma (x))
\]

belongs to \(\mathfrak{o}(\varPhi)\) . Put \(\tilde{H} = (-1)^{\sim}\) . Show that \((\tilde{H}, \tilde{X}_{+}, \tilde{X}_{-})\) is an \(\mathfrak{sl}_2\) -triplet in the Lie algebra \(\mathfrak{o}(\varPhi)\) .

\begin{enumerate}
\def\labelenumi{\alph{enumi})}
\setcounter{enumi}{2}
\tightlist
\item
  Show that \(\rho\) takes the endomorphisms \(\tilde{H}, \tilde{X}_{+}, \tilde{X}_{-}\) of \(\mathfrak{o}(\Phi)\) to the endomorphisms of N denoted by \(H, X_{+}\) and \(X_{-}\) , respectively, in no. 3.IV. Deduce that \((H, X_{+}, X_{-})\) is an \(\mathfrak{sl}_{2}\) -triplet in the Lie algebra \(\mathfrak{gl}(N)\) .
\end{enumerate}

\specialchapter{SUMMARY OF SOME IMPORTANT PROPERTIES OF SEMI-SIMPLE LIE ALGEBRAS}

In this summary, g denotes a semi-simple Lie algebra over k.

\section*{CARTAN SUBALGEBRAS}

\begin{enumerate}
\def\labelenumi{\arabic{enumi})}
\item
  Let E be the set of commutative subalgebras of g that are reductive in g; this is also the set of commutative subalgebras of g all of whose elements are semi-simple. The Cartan subalgebras of g are the maximal elements of E.
\item
  Let \(x\) be a regular element of \(\mathfrak{g}\) . Then \(x\) is semi-simple. There exists a unique Cartan subalgebra of \(\mathfrak{g}\) containing \(x\) ; it is the commutant of \(x\) in \(\mathfrak{g}\) .
\item
  Let \(x\) be a semi-simple element of \(\mathfrak{g}\) . Then \(x\) belongs to a Cartan subalgebra of \(\mathfrak{g}\) . The element \(x\) is regular if and only if the dimension of the commutant of \(x\) is equal to the rank of \(\mathfrak{g}\) .
\item
  Let \(\mathfrak{h}\) be a Cartan subalgebra of \(\mathfrak{g}\) . Then \(\mathfrak{h}\) is said to be splitting if \(\operatorname{ad}_{\mathfrak{g}}x\) is triangularizable for all \(x \in \mathfrak{h}\) . Moreover, \(\mathfrak{g}\) is said to be splittable if \(\mathfrak{g}\) has a splitting Cartan subalgebra (this is the case if \(k\) is algebraically closed). A split semi-simple Lie algebra is a pair \((\mathfrak{g}, \mathfrak{h})\) where \(\mathfrak{g}\) is a semi-simple Lie algebra and \(\mathfrak{h}\) is a splitting Cartan subalgebra of \(\mathfrak{g}\) .
\end{enumerate}

In the remainder of this summary, \((\mathfrak{g},\mathfrak{h})\) denotes a split semi-simple Lie algebra.

\section*{ROOT SYSTEMS}

\begin{enumerate}
\def\labelenumi{\arabic{enumi})}
\setcounter{enumi}{4}
\item
  For any element \(\alpha\) of the dual \(\mathfrak{h}^*\) of \(\mathfrak{h}\) , let \(\mathfrak{g}^{\alpha}\) be the set of \(x\in \mathfrak{g}\) such that \([h,x] = \alpha (h)x\) for all \(h\in \mathfrak{h}\) . If \(\alpha = 0\) , then \(\mathfrak{g}^{\alpha} = \mathfrak{h}\) . Any \(\alpha \in \mathfrak{h}^{*} - \{0\}\) such that \(\mathfrak{g}^{\alpha}\neq 0\) is called a root of \((\mathfrak{g},\mathfrak{h})\) . Denote by \(\mathrm{R}(\mathfrak{g},\mathfrak{h})\) (or simply by R) the set of roots of \((\mathfrak{g},\mathfrak{h})\) . This is a reduced root system in \(\mathfrak{h}^*\) in the sense of Chap. VI, §1, no. 4. The algebra \(\mathfrak{g}\) is simple if and only if R is irreducible.
\item
  For all \(\alpha \in \mathbb{R}\) , \(\mathfrak{g}^{\alpha}\) is of dimension 1. The vector space \([\mathfrak{g}^{\alpha}, \mathfrak{g}^{-\alpha}]\) is contained in \(\mathfrak{h}\) , is of dimension 1, and contains a unique element \(H_{\alpha}\) such that \(\alpha(H_{\alpha}) = 2\) ; we have \(H_{\alpha} = \alpha^{\vee}\) (Chap. VI, §1, no. 1); the set of \(H_{\alpha}\) , for \(\alpha \in \mathbb{R}\) , is the inverse root system \(\mathbb{R}^{\vee}\) of \(\mathbb{R}\) .
\item
  We have \(\mathfrak{g} = \mathfrak{h} \oplus \bigoplus_{\alpha \in \mathrm{R}} \mathfrak{g}^{\alpha}\) . There exists a family \((X_{\alpha})_{\alpha \in \mathrm{R}}\) such that, for all \(\alpha \in \mathrm{R}\) , we have \(X_{\alpha} \in \mathfrak{g}^{\alpha}\) and \([X_{\alpha}, X_{-\alpha}] = -H_{\alpha}\) . Every \(x \in \mathfrak{g}\) can be written uniquely in the form
\end{enumerate}

\[
x = h + \sum_ {\alpha \in \mathrm{R}} \lambda_ {\alpha} X _ {\alpha}, \quad \text { where } h \in \mathfrak {h}, \lambda_ {\alpha} \in k.
\]

The bracket of two elements can be calculated by means of the formulas

\[
\begin{array}{c} [ h, X _ {\alpha} ] = \alpha (h) X _ {\alpha} \\ \relax [ X _ {\alpha}, X _ {\beta} ] = 0 \quad \text {if} \alpha + \beta \notin \mathrm{R} \cup \{0 \} \\ \relax [ X _ {\alpha}, X _ {- \alpha} ] = - H _ {\alpha} \\ \relax [ X _ {\alpha}, X _ {\beta} ] = \mathrm{N} _ {\alpha \beta} X _ {\alpha + \beta} \quad \text {if} \alpha + \beta \in \mathrm{R}, \end{array}
\]

the \(\mathrm{N}_{\alpha \beta}\) being non-zero elements of \(k\) .

\begin{enumerate}
\def\labelenumi{\arabic{enumi})}
\setcounter{enumi}{7}
\tightlist
\item
  Let B be a basis of R. The algebra g is generated by the \(X_{\alpha}\) and the \(X_{-\alpha}\) for \(\alpha \in B\) . We have \([X_{\alpha}, X_{-\beta}] = 0\) if \(\alpha, \beta \in B\) and \(\alpha \neq \beta\) . Let \((n(\alpha, \beta))_{\alpha, \beta \in B}\) be the Cartan matrix of R (relative to B). We have \(n(\alpha, \beta) = \alpha(H_{\beta})\) . If \(\alpha, \beta \in B\) and \(\alpha \neq \beta\) , \(n(\alpha, \beta)\) is a negative integer and
\end{enumerate}

\[
(\operatorname{ad} X _ {\beta}) ^ {1 - n (\alpha , \beta)} X _ {\alpha} = 0 \quad \text { and } \quad (\operatorname{ad} X _ {- \beta}) ^ {1 - n (\alpha , \beta)} X _ {- \alpha} = 0.
\]

\begin{enumerate}
\def\labelenumi{\arabic{enumi})}
\setcounter{enumi}{8}
\tightlist
\item
  If \(\alpha, \beta, \alpha + \beta \in \mathbb{R}\) , let \(q_{\alpha \beta}\) be the largest integer \(j\) such that \(\beta - j\alpha \in \mathbb{R}\) . The family \((X_{\alpha})_{\alpha \in \mathbb{R}}\) in 7) can be chosen so that \(\mathrm{N}_{\alpha, \beta} = \mathrm{N}_{-\alpha, -\beta}\) if \(\alpha, \beta, \alpha + \beta \in \mathbb{R}\) . Then \(\mathrm{N}_{\alpha \beta} = \pm (q_{\alpha \beta} + 1)\) . There exists an involutive automorphism \(\theta\) of \(\mathfrak{g}\) that takes \(X_{\alpha}\) to \(X_{-\alpha}\) for all \(\alpha \in \mathbb{R}\) ; we have \(\theta(h) = -h\) for all \(h \in \mathfrak{h}\) . The \(\mathbf{Z}\) -submodule \(\mathfrak{g}_{\mathbf{Z}}\) of \(\mathfrak{g}\) generated by the \(H_{\alpha}\) and the \(X_{\alpha}\) is a \(\mathbf{Z}\) -Lie subalgebra of \(\mathfrak{g}\) , and the canonical map \(\mathfrak{g}_{\mathbf{Z}} \otimes_{\mathbf{Z}} k \to \mathfrak{g}\) is an isomorphism.
\end{enumerate}

The pair \((\mathfrak{g},\mathfrak{h})\) can be obtained by extension of scalars from a split semi-simple Q-Lie algebra.

\begin{enumerate}
\def\labelenumi{\arabic{enumi})}
\setcounter{enumi}{9}
\item
  The Weyl group, group of weights, \ldots{} of R is called the Weyl group, group of weights, \ldots{} of \((\mathfrak{g},\mathfrak{h})\) . The Weyl group will be denoted by W in what follows. We consider its operation, not only on \(h^{*}\) , but also on h (by transport of structure). If \(h_{Q}\) (resp. \(h_{Q}^{*}\) ) denotes the Q-vector subspace of h (resp. \(h^{*}\) ) generated by the \(H_{\alpha}\) (resp. the \(\alpha\) ), then h (resp. \(h^{*}\) ) can be canonically identified with \(h_{Q} \otimes_{Q} k\) (resp. \(h_{Q}^{*} \otimes_{Q} k\) ), and \(h_{Q}^{*}\) can be identified with the dual of \(h_{Q}\) . When we speak of the Weyl chambers of R, these are understood to be in \(h_{R} = h_{Q} \otimes_{Q} R\) or \(h_{R}^{*} = h_{Q}^{*} \otimes_{Q} R\) .
\item
  Let \(\varPhi\) be the Killing form of \(\mathfrak{g}\) . If \(\alpha + \beta \neq 0\) , \(\mathfrak{g}^{\alpha}\) and \(\mathfrak{g}^{\beta}\) are orthogonal with respect to \(\varPhi\) . The restriction of \(\varPhi\) to \(\mathfrak{g}^{\alpha} \times \mathfrak{g}^{-\alpha}\) is non-degenerate. If \(x, y \in \mathfrak{h}\) , then \(\varPhi(x, y) = \sum_{\alpha \in \mathrm{R}} \alpha(x)\alpha(y)\) . We have \(\varPhi(H_{\alpha}, H_{\beta}) \in \mathbf{Z}\) . The restriction of \(\varPhi\) to \(\mathfrak{h}\) is non-degenerate and invariant under W; its restriction to \(\mathfrak{h}_{\mathbf{Q}}\) is positive.
\item
  The root system of \((\mathfrak{g},\mathfrak{h})\) depends, up to isomorphism, only on g and not on h. By abuse of language, the Weyl group, group of weights, \ldots{} of \((\mathfrak{g},\mathfrak{h})\) are also called the Weyl group, group of weights, \ldots{} of g.
\end{enumerate}

If \(R_{1}\) is a reduced root system, there exists a split semi-simple Lie algebra \((\mathfrak{g}_{1},\mathfrak{h}_{1})\) such that \(\mathrm{R}(\mathfrak{g}_{1},\mathfrak{h}_{1})\) is isomorphic to \(R_{1}\) ; it is unique, up to isomorphism.

The classification of splittable semi-simple Lie algebras is thus reduced to that of root systems.

\section*{SUBALGEBRAS}

\begin{enumerate}
\def\labelenumi{\arabic{enumi})}
\setcounter{enumi}{12}
\item
  If \(P \subset R\) , put \(g^{P} = \bigoplus_{\alpha \in P} g^{\alpha}\) and \(h_{P} = \sum_{\alpha \in P} kH_{\alpha}\) . Let \(P \subset R\) , \(h'\) a vector subspace of h, and \(a = h' \oplus g^{P}\) . Then a is a subalgebra of g if and only if P is a closed subset of R and \(h'\) contains \(h_{P \cap (-P)}\) ; a is reductive in g if and only if P = -P; and a is solvable if and only if \(P \cap (-P) = \varnothing\) .
\item
  Let \(\mathbf{P}\) be a closed subset of \(\mathbf{R}\) , and \(\mathfrak{b} = \mathfrak{h} \oplus \mathfrak{g}^{\mathrm{P}}\) . The following conditions are equivalent:
\end{enumerate}

\begin{enumerate}
\def\labelenumi{(\roman{enumi})}
\item
  \(\mathfrak{b}\) is a maximal solvable subalgebra of \(\mathfrak{g}\) ;
\item
  \(\mathrm{P} \cap (-\mathrm{P}) = \varnothing\) and \(\mathrm{P} \cup (-\mathrm{P}) = \mathrm{R}\) ;
\item
  there exists a chamber \(\mathbf{C}\) of \(\mathbf{R}\) such that \(\mathrm{P} = \mathrm{R}_{+}(\mathrm{C})\) (cf.~Chap. VI, §1, no. 6).
\end{enumerate}

A Borel subalgebra of \((\mathfrak{g},\mathfrak{h})\) is a subalgebra of g containing h and satisfying the above conditions. A subalgebra b of g is called a Borel subalgebra of g if there exists a splitting Cartan subalgebra \(h'\) of g such that b is a Borel subalgebra of \((\mathfrak{g},\mathfrak{h}')\) ; if k is algebraically closed, this is equivalent to saying that b is a maximal solvable subalgebra of g.

Let \(\mathfrak{b} = \mathfrak{h} \oplus \mathfrak{g}^{\mathrm{R}_{+}(C)}\) be a Borel subalgebra of \((\mathfrak{g}, \mathfrak{h})\) . The largest nilpotent ideal of \(\mathfrak{b}\) is \([\mathfrak{b}, \mathfrak{b}] = \mathfrak{g}^{\mathrm{R}_{+}(C)}\) . Let \(B\) be the basis of \(R\) associated to \(C\) ; the algebra \([\mathfrak{b}, \mathfrak{b}]\) is generated by the \(\mathfrak{g}^{\alpha}\) for \(\alpha \in B\) .

If \(\mathfrak{b},\mathfrak{b}'\) are Borel subalgebras of \(\mathfrak{g}\) , there exists a Cartan subalgebra of \(\mathfrak{g}\) contained in \(\mathfrak{b} \cap \mathfrak{b}'\) ; such a subalgebra is splitting.

\begin{enumerate}
\def\labelenumi{\arabic{enumi})}
\setcounter{enumi}{14}
\tightlist
\item
  Let \(\mathbf{P}\) be a closed subset of \(\mathbf{R}\) , and \(\mathfrak{p} = \mathfrak{h} \oplus \mathfrak{g}^{\mathrm{P}}\) . The following conditions are equivalent:
\end{enumerate}

\begin{enumerate}
\def\labelenumi{(\roman{enumi})}
\item
  p contains a Borel subalgebra of \((\mathfrak{g}, \mathfrak{h})\) ;
\item
  \(P \cup (-P) = R;\)
\item
  there exists a chamber C of R such that \(\mathrm{P} \supset \mathrm{R}_{+}(\mathrm{C})\) .
\end{enumerate}

A parabolic subalgebra of \((\mathfrak{g},\mathfrak{h})\) is a subalgebra of g containing h and satisfying the above conditions. A subalgebra p of g is called parabolic if there exists a splitting Cartan subalgebra \(h'\) of g such that p is a parabolic subalgebra of \((\mathfrak{g},\mathfrak{h}')\) .

Let \(\mathfrak{p} = \mathfrak{h} \oplus \mathfrak{g}^{\mathrm{P}}\) be a parabolic subalgebra of \((\mathfrak{g}, \mathfrak{h})\) , Q the set of \(\alpha \in \mathrm{P}\) such that \(-\alpha \notin \mathrm{P}\) , and \(\mathfrak{s} = \mathfrak{h} \oplus \mathfrak{g}^{\mathrm{P} \cap (-\mathrm{P})}\) . Then \(\mathfrak{p} = \mathfrak{s} \oplus \mathfrak{g}^{\mathrm{Q}}\) , \(\mathfrak{s}\) is reductive in \(\mathfrak{g}\) , and \(\mathfrak{g}^{\mathrm{Q}}\) is the largest nilpotent ideal of \(\mathfrak{p}\) and the nilpotent radical of \(\mathfrak{p}\) . The centre of \(\mathfrak{p}\) is 0.

\section*{AUTOMORPHISMS}

\begin{enumerate}
\def\labelenumi{\arabic{enumi})}
\setcounter{enumi}{15}
\tightlist
\item
  The subgroup of \(\mathrm{Aut}(\mathfrak{g})\) generated by the \(e^{\mathrm{ad}x}\) , with \(x\) nilpotent, is the group \(\mathrm{Aut}_e(\mathfrak{g})\) of elementary automorphisms of \(\mathfrak{g}\) ; it is a normal subgroup of \(\mathrm{Aut}(\mathfrak{g})\) ; it is equal to its derived group.
\end{enumerate}

If \(\bar{k}\) is an algebraic closure of \(k\) , the group \(\operatorname{Aut}(\mathfrak{g})\) embeds naturally in \(\operatorname{Aut}(\mathfrak{g} \otimes_k \bar{k})\) . Put

\[
\operatorname{Aut} _ {0} (\mathfrak {g}) = \operatorname{Aut} (\mathfrak {g}) \cap \operatorname{Aut} _ {e} (\mathfrak {g} \otimes_ {k} \bar {k});
\]

this is a normal subgroup of \(\operatorname{Aut}(\mathfrak{g})\) , independent of the choice of \(\bar{k}\) . We have

\[
\operatorname{Aut} _ {e} (\mathfrak {g}) \subset \operatorname{Aut} _ {0} (\mathfrak {g}) \subset \operatorname{Aut} (\mathfrak {g}).
\]

The derived group of \(\mathrm{Aut}_0(\mathfrak{g})\) is \(\mathrm{Aut}_e(\mathfrak{g})\) . In the Zariski topology, \(\mathrm{Aut}(\mathfrak{g})\) and \(\mathrm{Aut}_0(\mathfrak{g})\) are closed in \(\mathrm{End}_k(\mathfrak{g})\) , \(\mathrm{Aut}_0(\mathfrak{g})\) is the connected component of the identity in \(\mathrm{Aut}(\mathfrak{g})\) , and \(\mathrm{Aut}_e(\mathfrak{g})\) is dense in \(\mathrm{Aut}_0(\mathfrak{g})\) .

Let B be a basis of R, and \(\operatorname{Aut}(R,B)\) the group of automorphisms of R that leave B stable. Then \(\operatorname{Aut}(\mathfrak{g})\) is the semi-direct product of a subgroup isomorphic to \(\operatorname{Aut}(R,B)\) and \(\operatorname{Aut}_{0}(\mathfrak{g})\) ; in particular, \(\operatorname{Aut}(\mathfrak{g})/\operatorname{Aut}_{0}(\mathfrak{g})\) is isomorphic to \(\operatorname{Aut}(R,B)\) , which is itself isomorphic to a group of automorphisms of the Dynkin graph of g.

\begin{enumerate}
\def\labelenumi{\arabic{enumi})}
\setcounter{enumi}{16}
\tightlist
\item
  A framing of g is a triplet \((\mathfrak{h}^{\prime},\mathrm{B},(X_{\alpha})_{\alpha\in\mathrm{B}})\) , where \(h^{\prime}\) is a splitting Cartan subalgebra of g, B is a basis of \(\mathrm{R}(\mathfrak{g},\mathfrak{h}^{\prime})\) , and where, for all \(\alpha\in B\) , \(X_{\alpha}\) is a non-zero element of \(g^{\alpha}\) . The group \(\operatorname{Aut}_{0}(\mathfrak{g})\) operates simply-transitively on the set of framings of g.
\end{enumerate}

The group \(\mathrm{Aut}_e(\mathfrak{g})\) operates transitively on the set of pairs \((\mathfrak{k},\mathfrak{b})\) , where \(\mathfrak{k}\) is a splitting Cartan subalgebra of \(\mathfrak{g}\) and \(\mathfrak{b}\) is a Borel subalgebra of \((\mathfrak{g},\mathfrak{k})\) .

\begin{enumerate}
\def\labelenumi{\arabic{enumi})}
\setcounter{enumi}{17}
\tightlist
\item
  Denote by \(\operatorname{Aut}(\mathfrak{g},\mathfrak{h})\) the set of \(s\in \operatorname{Aut}(\mathfrak{g})\) such that \(s(\mathfrak{h}) = \mathfrak{h}\) . Put
\end{enumerate}

\[
\operatorname{Aut} (\mathfrak {g}, \mathfrak {h}) = \operatorname{Aut} _ {e} (\mathfrak {g}) \cap \operatorname{Aut} (\mathfrak {g}, \mathfrak {h}), \quad \operatorname{Aut} _ {0} (\mathfrak {g}, \mathfrak {h}) = \operatorname{Aut} _ {0} (\mathfrak {g}) \cap \operatorname{Aut} (\mathfrak {g}, \mathfrak {h}).
\]

If \(s \in \text{Aut}(\mathfrak{g}, \mathfrak{h})\) , the contragredient map of \(s | \mathfrak{h}\) is an element of the group \(\text{A(R)}\) of automorphisms of R; denote this element by \(\varepsilon(s)\) ; the map \(\varepsilon\) is a homomorphism from \(\text{Aut}(\mathfrak{g}, \mathfrak{h})\) to \(\text{A(R)}\) . We have \(\text{Aut}_{0}(\mathfrak{g}) = \text{Aut}_{e}(\mathfrak{g})\) . Ker \(\varepsilon\) , and

\[
\varepsilon (\operatorname{Aut} _ {0} (\mathfrak {g}, \mathfrak {h})) = \varepsilon (\operatorname{Aut} _ {e} (\mathfrak {g}, \mathfrak {h})) = W.
\]

Let \(\mathrm{T}_{\mathrm{P}} = \mathrm{Hom}(\mathrm{P}(\mathrm{R}), k^{*})\) , \(\mathrm{T}_{\mathrm{Q}} = \mathrm{Hom}(\mathrm{Q}(\mathrm{R}), k^{*})\) . The injection of \(\mathrm{Q}(\mathrm{R})\) into \(\mathrm{P}(\mathrm{R})\) defines a homomorphism from \(\mathrm{T}_{\mathrm{P}}\) to \(\mathrm{T}_{\mathrm{Q}}\) ; let \(\mathrm{Im}(\mathrm{T}_{\mathrm{P}})\) be its image. If \(t \in \mathrm{T}_{\mathrm{Q}}\) , let \(f(t)\) be the endomorphism of \(\mathfrak{g}\) such that, for all \(\alpha \in \mathrm{R} \cup \{0\}\) , \(f(t)|\mathfrak{g}^{\alpha}\) is the homothety with ratio \(t(\alpha)\) ; we have \(f(t)\in \mathrm{Aut}_{0}(\mathfrak{g},\mathfrak{h})\) and \(f\) is an injective homomorphism from \(\mathrm{T}_{\mathrm{Q}}\) to \(\mathrm{Aut}_0(\mathfrak{g},\mathfrak{h})\) . The sequences

\[
\{1 \} \longrightarrow \mathrm {T_ {Q}} \stackrel {{f}} {{\longrightarrow}} \operatorname{Aut} (\mathfrak {g}, \mathfrak {h}) \stackrel {{\varepsilon}} {{\longrightarrow}} \mathrm{A(R)} \longrightarrow \{1 \}
\]

and

\[
\{1 \} \longrightarrow \mathrm{T} _ {\mathrm{Q}} \stackrel {f} {\longrightarrow} \mathrm{Aut} _ {0} (\mathfrak {g}, \mathfrak {h}) \stackrel {\varepsilon} {\longrightarrow} \mathrm{W} \longrightarrow \{1 \}
\]

are exact. We have \(f(\mathrm{Im}(\mathrm{T}_{\mathrm{P}})) \subset \mathrm{Aut}_{e}(\mathfrak{g}, \mathfrak{h})\) ; f defines, by passage to the quotient, a surjective \(^{23}\) homomorphism \(\mathrm{T}_{\mathrm{Q}}/\mathrm{Im}(\mathrm{T}_{\mathrm{P}}) \to \mathrm{Aut}_{0}(\mathfrak{g})/\mathrm{Aut}_{e}(\mathfrak{g})\) . In the Zariski topology, \(f(\mathrm{T}_{\mathrm{Q}})\) is closed in \(\mathrm{Aut}(\mathfrak{g})\) , and \(f(\mathrm{Im}(\mathrm{T}_{\mathrm{P}}))\) is dense in \(f(\mathrm{T}_{\mathrm{Q}})\) .

\section*{FINITE DIMENSIONAL MODULES}

\begin{enumerate}
\def\labelenumi{\arabic{enumi})}
\setcounter{enumi}{18}
\item
  Let V be a finite dimensional g-module. For all \(\mu\in h^{*}\) , let \(V^{\mu}\) be the set of \(v\in V\) such that \(h.v=\mu(h)v\) for all \(h\in h\) . The dimension of \(V^{\mu}\) is called the multiplicity of \(\mu\) in V; if it is \(\geq1\) , i.e.~if \(V^{\mu}\neq0\) , \(\mu\) is said to be a weight of V. We have \(V=\bigoplus_{\mu\in h^{*}}V^{\mu}\) . Every weight of V belongs to P(R). If \(\mu\) is a weight of V, and if \(w\in W\) , \(w\mu\) is a weight of V of the same multiplicity as \(\mu\) . If \(v\in V^{\mu}\) and \(x\in g^{\alpha}\) , then \(x.v\in V^{\mu+\alpha}\) .
\item
  Let B be a basis of R. Giving B determines an order relation on \(h_{Q}^{*}\) : the elements of \(h_{Q}^{*}\) that are \(\geq 0\) are the linear combinations of elements of B with rational coefficients \(\geq 0\) . Denote by \(Q_{+}(R)\) (resp. \(R_{+}\) ) the set of positive elements of \(Q(R)\) (resp. of R).
\end{enumerate}

Let \(V\) be a finite dimensional simple \(\mathfrak{g}\) -module. Then \(V\) has a highest weight \(\lambda\) . This weight is of multiplicity 1, and it is a dominant weight: if \(\alpha \in R_{+}\) , \(\lambda(H_{\alpha})\) is an integer \(\geq 0\) . We have \(\mathfrak{g}^{\alpha}V^{\lambda} = 0\) if \(\alpha \in R_{+}\) . Every weight of \(V\) is of the form \(\lambda - \nu\) with \(\nu \in Q_{+}(R)\) ; conversely, if a weight is dominant and is of the form \(\lambda - \nu\) with \(\nu \in Q_{+}(R)\) , then it is a weight of \(V\) .

\begin{enumerate}
\def\labelenumi{\arabic{enumi})}
\setcounter{enumi}{20}
\tightlist
\item
  Two finite dimensional simple g-modules with the same highest weight are isomorphic. Every dominant weight is the highest weight of a finite dimensional simple g-module.
\end{enumerate}

Every finite dimensional simple g-module is absolutely simple.

\begin{enumerate}
\def\labelenumi{\arabic{enumi})}
\setcounter{enumi}{21}
\item
  Let \(\varPhi\) be the Killing form of g, \(C \in U(\mathfrak{g})\) the corresponding Casimir element, \(\langle\cdot,\cdot\rangle\) the inverse form on \(h^{*}\) of \(\varPhi|h \times h\) , and \(\rho = \frac{1}{2} \sum_{\alpha \in R_{+}} \alpha\) . Let V be a finite dimensional simple g-module, of highest weight \(\lambda\) . Then \(C_{V}\) is the homothety with ratio \(\langle\lambda, \lambda + 2\rho\rangle\) .
\item
  Let \(\mathrm{V}\) be a finite dimensional \(\mathfrak{g}\) -module, and \(\mathrm{V}^*\) its dual. Then \(\mu \in \mathfrak{h}^*\) is a weight of \(\mathrm{V}^*\) if and only if \(-\mu\) is a weight of \(\mathrm{V}\) , and the multiplicity of \(\mu \in \mathrm{V}^*\) is equal to the multiplicity of \(-\mu\) in V. If V is simple of highest weight \(\lambda\) , \(V^{*}\) is simple of highest weight \(-w_{0}\lambda\) , where \(w_{0}\) is the element of W that takes B to -B.
\item
  Let \(\mathrm{V}\) be a finite dimensional simple \(\mathfrak{g}\) -module of highest weight \(\lambda\) , and \(\mathcal{B}\) the vector space of \(\mathfrak{g}\) -invariant bilinear forms on \(\mathrm{V}\) . Let \(m\) be the integer \(\sum_{\alpha \in \mathbb{R}_+} \lambda(H_\alpha)\) , and let \(w_0 \in \mathrm{W}\) be as in 23). If \(w_0 \lambda \neq -\lambda\) , \(\mathrm{V}\) and \(V^*\) are not isomorphic, and \(\mathcal{B} = 0\) . If \(w_0 \lambda = -\lambda\) , then \(\dim \mathcal{B} = 1\) and every non-zero element of \(\mathcal{B}\) is non-degenerate; if \(m\) is even (resp. odd), every element of \(\mathcal{B}\) is symmetric (resp. alternating).
\item
  Let \(\mathbf{Z}[\mathrm{P}]\) be the algebra of the group \(\mathrm{P} = \mathrm{P}(\mathrm{R})\) with coefficients in \(\mathbf{Z}\) . If \(\lambda \in \mathrm{P}\) , denote by \(e^{\lambda}\) the corresponding element of \(\mathbf{Z}[\mathrm{P}]\) ; the \(e^{\lambda}\) , \(\lambda \in \mathrm{P}\) , form a \(\mathbf{Z}\) -basis of \(\mathbf{Z}[\mathrm{P}]\) , and \(e^{\lambda}e^{\mu} = e^{\lambda +\mu}\) for \(\lambda, \mu \in \mathrm{P}\) .
\end{enumerate}

Let \(\mathrm{V}\) be a finite dimensional \(\mathfrak{g}\) -module. The character of \(\mathrm{V}\) , denoted by \(\operatorname{ch} \mathrm{V}\) , is the element \(\sum_{\mu \in \mathrm{P}} (\dim \mathrm{V}^{\mu}) e^{\mu}\) of \(\mathbf{Z}[\mathrm{P}]\) ; this element belongs to the subalgebra \(\mathbf{Z}[\mathrm{P}]^{\mathrm{W}}\) of \(\mathbf{Z}[\mathrm{P}]\) consisting of the elements invariant under \(\mathrm{W}\) . We have

\[
\operatorname{ch} \left(\mathrm{V} \oplus \mathrm{V} ^ {\prime}\right) = \operatorname{ch} \mathrm{V} + \operatorname{ch} \mathrm{V} ^ {\prime} \quad \text { and } \quad \operatorname{ch} \left(\mathrm{V} \otimes \mathrm{V} ^ {\prime}\right) = (\operatorname{ch} \mathrm{V}). (\operatorname{ch} \mathrm{V} ^ {\prime}).
\]

Two finite dimensional g-modules with the same character are isomorphic.

For all \(\alpha\in B\) , let \(V_{\alpha}\) be a simple g-module with highest weight the fundamental weight \(\varpi_{\alpha}\) corresponding to \(\alpha\) . The elements ch \(V_{\alpha},\alpha\in B\) , are algebraically independent and generate the Z-algebra \(Z[P]^{W}\) .

\begin{enumerate}
\def\labelenumi{\arabic{enumi})}
\setcounter{enumi}{25}
\tightlist
\item
  Let \(\rho\) be half the sum of the roots \(\geq 0\) . For all \(w \in W\) , let \(\varepsilon(w)\) be the determinant of w, equal to \(\pm 1\) . If V is a finite dimensional simple g-module of highest weight \(\lambda\) , then
\end{enumerate}

\[
\left(\sum_ {w \in \mathrm{W}} \varepsilon (w) e ^ {w \rho}\right). \mathrm{chV} = \sum_ {w \in \mathrm{W}} \varepsilon (w) e ^ {w (\lambda + \rho)},
\]

and

\[
\dim V = \prod_ {\alpha \in R _ {+}} \frac {\langle \lambda + \rho , H _ {\alpha} \rangle}{\langle \rho , H _ {\alpha} \rangle}.
\]

\begin{enumerate}
\def\labelenumi{\arabic{enumi})}
\setcounter{enumi}{26}
\tightlist
\item
  For all \(\nu \in \mathrm{P}\) , let \(\mathfrak{P}(\nu)\) be the number of families \((n_{\alpha})_{\alpha \in \mathrm{R}_{+}}\) , where the \(n_{\alpha}\) are integers \(\geq 0\) such that \(\nu = \sum_{\alpha \in \mathrm{R}_{+}} n_{\alpha} \alpha\) . Let \(\mathrm{V}\) be a finite dimensional simple \(\mathfrak{g}\) -module of highest weight \(\lambda\) . If \(\mu \in \mathrm{P}\) , the multiplicity of \(\mu\) in \(\mathrm{V}\) is
\end{enumerate}

\[
\sum_ {w \in \mathrm{W}} \varepsilon (w) \mathfrak {P} (w (\lambda + \rho) - (\mu + \rho)).
\]

\begin{enumerate}
\def\labelenumi{\arabic{enumi})}
\setcounter{enumi}{27}
\tightlist
\item
  Let \(\mathrm{V}, \mathrm{V}', \mathrm{V}''\) be finite dimensional simple \(\mathfrak{g}\) -modules, \(\lambda, \mu, \nu\) their highest weights. In \(\mathrm{V} \otimes \mathrm{V}'\) , the isotypical component of type \(\mathrm{V}''\) has length
\end{enumerate}

\[
\sum_ {w, w ^ {\prime} \in \mathrm{W}} \varepsilon (w w ^ {\prime}) \mathfrak {P} (w (\lambda + \rho) + w ^ {\prime} (\mu + \rho) - (\nu + 2 \rho)).
\]

In particular, if \(\nu = \lambda + \mu\) , the isotypical component in question is simple, and generated by \((\mathrm{V} \otimes \mathrm{V}')^{\lambda + \mu} = \mathrm{V}^{\lambda} \otimes \mathrm{V}'^{\mu}\) .

\section*{INVARIANT POLYNOMIAL FUNCTIONS}

\begin{enumerate}
\def\labelenumi{\arabic{enumi})}
\setcounter{enumi}{28}
\item
  The algebra of polynomial functions on \(\mathfrak{g}\) can be identified with the symmetric algebra \(\mathbf{S}(\mathfrak{g}^*)\) of \(\mathfrak{g}^*\) , and hence is a \(\mathfrak{g}\) -module in a canonical way; hence the notion of an invariant polynomial function on \(\mathfrak{g}\) . Let \(f \in \mathbf{S}(\mathfrak{g}^*)\) . Then \(f\) is invariant if and only if \(f \circ s = f\) for all \(s \in \mathrm{Aut}_0(\mathfrak{g})\) , or that \(f \circ s = f\) for all \(s \in \mathrm{Aut}_e(\mathfrak{g})\) .
\item
  Let \(\mathrm{I}(\mathfrak{g}^*)\) be the algebra of invariant polynomial functions on \(\mathfrak{g}\) , and \(\mathbf{S}(\mathfrak{h}^*)^{\mathrm{W}}\) the algebra of W-invariant polynomial functions on \(\mathfrak{h}\) . Let
\end{enumerate}

\[
i: \mathbf {S} (\mathfrak {g} ^ {*}) \to \mathbf {S} (\mathfrak {h} ^ {*})
\]

be the restriction homomorphism. The map \(i|I(\mathfrak{g}^{*})\) is an isomorphism from \(I(\mathfrak{g}^{*})\) to \(\mathbf{S}(\mathfrak{h}^{*})^{\mathrm{W}}\) . If l is the rank of g, there exist l homogeneous elements of \(I(\mathfrak{g}^{*})\) that are algebraically independent and generate the algebra \(I(\mathfrak{g}^{*})\) .

\begin{enumerate}
\def\labelenumi{\arabic{enumi})}
\setcounter{enumi}{30}
\item
  An element \(a\) of \(\mathfrak{g}\) is nilpotent if and only if \(f(a) = 0\) for every homogeneous element \(f\) of \(\mathrm{I}(\mathfrak{g}^*)\) of degree \(> 0\) .
\item
  Let \(s \in \mathrm{Aut}(\mathfrak{g})\) . Then \(s\) belongs to \(\mathrm{Aut}_0(\mathfrak{g})\) if and only if \(f \circ s = f\) for all \(f \in \mathrm{I}(\mathfrak{g}^*)\) .
\end{enumerate}

\section*{\texorpdfstring{\(\mathfrak{sl}_2\) -TRIPLETS}{\textbackslash mathfrak\{sl\}\_2 -TRIPLETS}}

\begin{enumerate}
\def\labelenumi{\arabic{enumi})}
\setcounter{enumi}{32}
\item
  An \(sl_{2}\) -triplet in g is a sequence \((x,h,y)\) of elements of g distinct from \((0,0,0)\) and such that \([h,x]=2x,[h,y]=-2y,[x,y]=-h\) . Then x,y are nilpotent in g, and h is semi-simple in g.
\item
  Let \(x\) be a non-zero nilpotent element of \(\mathfrak{g}\) . There exist \(h, y \in \mathfrak{g}\) such that \((x, h, y)\) is an \(\mathfrak{sl}_2\) -triplet.
\item
  Let \((x,h,y)\) and \((x',h',y')\) be \(\mathfrak{sl}_2\) -triplets in \(\mathfrak{g}\) . The following conditions are equivalent:
\end{enumerate}

\begin{enumerate}
\def\labelenumi{\alph{enumi})}
\item
  there exists \(s \in \text{Aut}_{e}(\mathfrak{g})\) such that \(sx = x'\) ;
\item
  there exists \(s \in \mathrm{Aut}_e(\mathfrak{g})\) such that \(sx = x', sh = h', sy = y'\) .
\end{enumerate}

\begin{enumerate}
\def\labelenumi{\arabic{enumi})}
\setcounter{enumi}{35}
\tightlist
\item
  If \(k\) is algebraically closed, conditions \(a)\) and \(b)\) of 35) are equivalent to: \(c)\) there exists \(s \in \operatorname{Aut}_e(\mathfrak{g})\) such that \(sh = h'\) .
\end{enumerate}

Moreover, the number of conjugacy classes, relative to \(\operatorname{Aut}_{e}(\mathfrak{g})\) , of non-zero nilpotent elements of g is at most equal to \(3^{l}\) , where l is the rank of g.

\begin{enumerate}
\def\labelenumi{\arabic{enumi})}
\setcounter{enumi}{36}
\tightlist
\item
  A nilpotent element x of g is called principal if the dimension of the centralizer of x is equal to the rank of g. There exist principal nilpotent elements in g. If k is algebraically closed, the principal nilpotent elements of g are conjugate under \(\operatorname{Aut}_{e}(\mathfrak{g})\) .
\end{enumerate}

\chapter{Compact Real Lie Groups}

In this chapter \(^{1}\) , the expression ``Lie group'' means ``finite dimensional Lie group over the field of real numbers'', the expression ``Lie algebra'' means, unless stated otherwise, ``finite dimensional Lie algebra over the field of real numbers'', the expression ``real Lie algebra'' (resp. ``complex Lie algebra'') means ``finite dimensional Lie algebra over the field of real numbers (resp. ``finite dimensional Lie algebra over the field of complex numbers'').

We denote by \(G_{0}\) the identity component of a topological group G. We denote by \(\mathrm{C(G)}\) the centre of a group G, by \(\mathrm{D(G)}\) its derived group, and by \(\mathrm{N_{G}(H)}\) or \(\mathrm{N(H)}\) (resp. \(\mathrm{Z_{G}(H)}\) or \(\mathrm{Z(H)}\) ) the normalizer (resp. centralizer) of a subset H of a group G.

\section{COMPACT LIE ALGEBRAS}

\subsection{INVARIANT HERMITIAN FORMS}

In this number, the letter k denotes the field R or C. Let V be a finite dimensional k-vector space, \(\Phi\) a separating \(^{2}\) positive hermitian form on V, G a group, g an R-Lie algebra, \(\rho: G \to \mathbf{GL}(V)\) a group homomorphism, \(\varphi: \mathfrak{g} \to \mathfrak{gl}(V)\) a homomorphism of R-Lie algebras.

\begin{enumerate}
\def\labelenumi{\alph{enumi})}
\tightlist
\item
  The form \(\Phi\) is invariant under G (resp. g) if and only if \(\rho(g)\) is unitary with respect to \(\Phi\) for all \(g \in G\) (resp. \(\varphi(x)\) is anti-hermitian \(^{3}\) with respect to \(\Phi\) for all \(x \in g\) ). Indeed, denote by \(a^{*}\) the adjoint of an endomorphism a of V with respect to \(\Phi\) ; for g in G, x in g, u and v in V, we have
\end{enumerate}

\[
\begin{array}{c} \Phi (\rho (g) u, \rho (g) v) = \Phi (\rho (g) ^ {*} \rho (g) u, v), \\ \Phi (\varphi (x) u, v) + \Phi (u, \varphi (x) v) = \Phi ((\varphi (x) + \varphi (x) ^ {*}). u, v); \end{array}
\]

thus, \(\Phi(\rho(g)u,\rho(g)v)=\Phi(u,v)\) for all \(u,v\) in V if and only if \(\rho(g)^{*}\rho(g)=\mathrm{Id}_{\mathrm{V}}\) ; similarly, \(\Phi(\varphi(x)u,v)+\Phi(u,\varphi(x)v)=0\) for all \(u,v\) in V if and only if \(\varphi(x)+\varphi(x)^{*}=0\) , hence the stated assertion.

\begin{enumerate}
\def\labelenumi{\alph{enumi})}
\setcounter{enumi}{1}
\item
  If the form \(\varPhi\) is invariant under G (resp. g), the orthogonal complement of a stable subspace of V is stable; in particular, the representation \(\rho\) (resp. \(\varphi\) ) is then semi-simple (cf.~Algebra, Chap. IX); moreover, for all \(g \in G\) (resp. \(x \in g\) ), the endomorphism \(\rho(g)\) (resp. \(\varphi(x)\) ) of V is then semi-simple, with eigenvalues of absolute value 1 (resp. with purely imaginary eigenvalues); indeed \(\rho(g)\) is unitary (resp. \(i\varphi(x)\) is hermitian, cf.~Algebra, Chap. IX).
\item
  Assume that \(k = \mathbf{R}\) . If \(G\) is a connected Lie group, \(\rho\) a morphism of Lie groups, \(\mathfrak{g}\) the Lie algebra of \(G\) and \(\varphi\) the homomorphism induced by \(\rho\) , then \(\varPhi\) is invariant under \(G\) if and only if it is invariant under \(\mathfrak{g}\) (Chap. III, §6, no. 5, Cor. 3).
\item
  There exists a separating positive hermitian form on V invariant under G if and only if the subgroup \(\rho(\mathrm{G})\) of \(\mathbf{GL}(\mathrm{V})\) is relatively compact (Integration, Chap. VII, §3, no. 1, Prop. 1).
\end{enumerate}

\subsection{CONNECTED COMMUTATIVE REAL LIE GROUPS}

Let \(G\) be a connected commutative (real) Lie group. The exponential map

\[
\exp_ {\mathrm{G}}: \mathrm{L} (\mathrm{G}) \to \mathrm{G}
\]

is a morphism of Lie groups, surjective with discrete kernel (Chap. III, §6, no. 4, Prop. 11), hence the fact that \(\mathrm{L}(\mathrm{G})\) is a connected covering of \(\mathrm{G}\) .

\begin{enumerate}
\def\labelenumi{\alph{enumi})}
\item
  The following conditions are equivalent: G is simply-connected, \(\exp_{G}\) is an isomorphism, G is isomorphic to \(R^{n}\) ( \(n = \dim G\) ). In this case, transporting the vector space structure of L(G) to G by the isomorphism \(\exp_{G}\) gives a vector space structure on G, which is the only one compatible with the topological group structure of G. Simply-connected commutative Lie groups are called vector (Lie) groups; unless stated otherwise, they are always given the R-vector space structure defined above.
\item
  Denote by \(\Gamma(\mathrm{G})\) the kernel of \(\exp_{G}\) . By General Topology, Chap. VII, §1, no. 1, Th. 1, the group G is compact if and only if \(\Gamma(\mathrm{G})\) is a lattice in \(\mathrm{L}(\mathrm{G})\) , in other words (loc. cit.) if the rank of the free Z-module \(\Gamma(\mathrm{G})\) is equal to the dimension of G. Conversely, if L is a finite dimensional R-vector space and \(\Gamma\) a lattice in L, the quotient topological group \(L/\Gamma\) is a compact connected commutative Lie group.
\end{enumerate}

The compact connected commutative Lie groups are called real tori, or (in this chapter) tori.

\begin{enumerate}
\def\labelenumi{\alph{enumi})}
\setcounter{enumi}{2}
\tightlist
\item
  In the general case, let E be the vector subspace of L(G) generated by \(\Gamma(\mathrm{G})\) , and let V be a complementary subspace. Then G is the direct product of its Lie subgroups \(\exp(\mathrm{E})\) and \(\exp(\mathrm{V})\) ; the first is a torus, the second is vector. Finally, every compact subgroup of G is contained in \(\exp(\mathrm{E})\) (since its projection onto \(\exp(V)\) is necessarily reduced to the identity element); thus, the subgroup \(\exp(E)\) is the unique maximal compact subgroup of G.
\end{enumerate}

For example, take \(\mathrm{G} = \mathrm{C}^*\) ; identify \(\mathrm{L(G)}\) with \(\mathbf{C}\) so that the exponential map of \(\mathrm{G}\) is \(x \mapsto e^x\) . Then \(\Gamma(\mathrm{G}) = 2\pi i\mathbf{Z}, \mathrm{E} = i\mathbf{R}\) , and so \(\exp(\mathrm{E}) = \mathbf{U}\) ; if we take \(\mathrm{V} = \mathbf{R}\) , then \(\exp(\mathrm{V}) = \mathbf{R}_+^*\) and we recover the isomorphism \(\mathbf{C}^* \to \mathbf{U} \times \mathbf{R}_+^*\) constructed in General Topology, Chap. VIII, §1, no. 3, Prop. 1.

\(d\) ) Note finally that \(\exp_{\mathrm{G}}:\mathrm{L}(\mathrm{G})\to \mathrm{G}\) is a universal covering of \(\mathbf{G}\) , hence \(\Gamma (\mathbf{G})\) can be identified naturally with the fundamental group of \(\mathbf{G}\) .

\subsection{COMPACT LIE ALGEBRAS}

PROPOSITION 1. Let \(\mathfrak{g}\) be a (real) Lie algebra. The following conditions are equivalent:

\begin{enumerate}
\def\labelenumi{(\roman{enumi})}
\item
  g is isomorphic to the Lie algebra of a compact Lie group.
\item
  The group \(\operatorname{Int}(\mathfrak{g})\) (Chap. III, §6, no. 2, Def. 2) is compact.
\item
  \(\mathfrak{g}\) has an invariant bilinear form (Chap. I, §3, no. 6) that is symmetric, positive and separating.
\item
  g is reductive (Chap. I, §6, no. 4, Def. 4); for all \(x \in g\) , the endomorphism ad x is semi-simple, with purely imaginary eigenvalues.
\item
  \(\mathfrak{g}\) is reductive and its Killing form B is negative.
\item
  \(\Longrightarrow\) (ii): if g is the Lie algebra of a compact Lie group G, the group \(\operatorname{Int}(\mathfrak{g})\) is separating and isomorphic to a quotient of the compact group \(G_{0}\) (Chap. III, §6, no. 4, Cor. 4), hence is compact.
\item
  \(\Longrightarrow\) (iii): if the group \(\operatorname{Int}(\mathfrak{g})\) is compact, there exists a symmetric bilinear form on g that is positive, separating and invariant under \(\operatorname{Int}(\mathfrak{g})\) (no. 1), hence also invariant under the adjoint representation of g.
\item
  \(\Longrightarrow\) (iv): if (iii) is satisfied, the adjoint representation of g is semi-simple (no. 1), hence g is reductive; moreover, the endomorphisms ad x, for \(x \in g\) , have the indicated properties (no. 1).
\item
  \(\Longrightarrow\) (v): for all \(x\in \mathfrak{g}\) , \(\mathrm{B}(x,x) = \mathrm{Tr}((\mathrm{ad}x)^2)\) ; consequently, \(\mathrm{B}(x,x)\) is the sum of the squares of the eigenvalues of ad \(x\) , and hence is negative if these are purely imaginary.
\item
  \(\Longrightarrow\) (i): assume that g is reductive, hence the product of a commutative subalgebra c and a semi-simple subalgebra s (Chap. I, §6, no. 4, Prop. 5). The Killing form of s is the restriction of the form B to s, hence is negative and separating if B is negative. The subgroup \(\operatorname{Int}(\mathfrak{s})\) of \(\mathbf{GL}(\mathfrak{s})\) is closed (it is the identity component of \(\operatorname{Aut}(\mathfrak{s})\) , Chap. III, §10, no. 2, Cor. 2) and leaves the separating positive form -B invariant; thus, it is compact, and s is isomorphic to the Lie algebra of the compact Lie group \(\operatorname{Int}(\mathfrak{s})\) . Further, since c is commutative, it is isomorphic to the Lie algebra of a torus T. Thus g is isomorphic to the Lie algebra of the compact Lie group \(\operatorname{Int}(\mathfrak{s}) \times \operatorname{T}\) .
\end{enumerate}

DEFINITION 1. A compact Lie algebra \(^{4}\) is a Lie algebra that has properties (i) to (v) of Proposition 1.

Thus, the compact Lie algebras are the products of a commutative algebra with a compact semi-simple algebra. In other words, a Lie algebra is compact if and only if it is reductive and its derived Lie algebra is compact.

The Lie algebra of a compact Lie group is compact.

PROPOSITION 2. a) The product of a finite number of Lie algebras is a compact Lie algebra if and only if each factor is compact.

\begin{enumerate}
\def\labelenumi{\alph{enumi})}
\setcounter{enumi}{1}
\item
  A subalgebra of a compact Lie algebra is compact.
\item
  Let \(\mathfrak{h}\) be an ideal of a compact Lie algebra \(\mathfrak{g}\) . Then the algebra \(\mathfrak{g} / \mathfrak{h}\) is compact and the extension \(\mathfrak{h} \to \mathfrak{g} \to \mathfrak{g} / \mathfrak{h}\) is trivial.
\end{enumerate}

Assertions \(a\) and \(b\) follow from the characterization (iii) of Prop. 1. Part \(c\) follows from \(a\) and the fact that, in a reductive Lie algebra, every ideal is a direct factor (Chap. I, §6, no. 4, Cor. of Prop. 5).

PROPOSITION 3. Let G be a Lie group of which the group of connected components is finite. The following conditions are equivalent:

\begin{enumerate}
\def\labelenumi{(\roman{enumi})}
\item
  The Lie algebra \(\mathrm{L}(\mathrm{G})\) is compact.
\item
  The group \(\mathrm{Ad}(\mathrm{G})\) is compact.
\item
  There exists a separating positive symmetric bilinear form on L(G) invariant under the adjoint representation of G.
\end{enumerate}

*(iv) G has a riemannian metric invariant under left and right translations.*

\begin{enumerate}
\def\labelenumi{(\roman{enumi})}
\item
  \(\Longrightarrow\) (ii): if L(G) is compact, the group \(\mathrm{Ad}(\mathrm{G}_{0}) = \mathrm{Int}(\mathrm{L}(\mathrm{G}))\) is compact; since it has finite index in \(\mathrm{Ad}(\mathrm{G})\) , this latter group is also compact.
\item
  \(\Longrightarrow\) (iii): this follows from no. 1.
\item
  \(\Longrightarrow\) (i): since \(\operatorname{Int}(\mathrm{L}(\mathrm{G})) \subset \operatorname{Ad}(\mathrm{G})\) , this follows from the characterization (iii) of Prop. 1.
\end{enumerate}

\(^{*}\) (iii) \(\Longleftrightarrow\) (iv): this follows from Chap. III, §3, no. 13.*

\subsection{GROUPS WHOSE LIE ALGEBRA IS COMPACT}

THEOREM 1. (H. Weyl) Let G be a connected Lie group whose Lie algebra is compact semi-simple. Then G is compact and its centre is finite.

Since G is semi-simple, its centre D is discrete. Moreover, the quotient group G/D is isomorphic to \(\operatorname{Ad}(G)\) (Chap. III, §6, no. 4, Cor. 4), hence compact (Prop. 3). Finally, the group G/D is equal to its derived group (Chap. III, §9, no. 2, Cor. of Prop. 4). The theorem now follows from Integration, Chap. VII, §3, no. 2, Prop. 5.

PROPOSITION 4. Let G be a connected Lie group whose Lie algebra is compact. There exist a torus T, a simply-connected compact semi-simple Lie group S, a vector group V and a surjective morphism \(f: V \times T \times S \to G\) with finite kernel. If G is compact, the group V is reduced to the identity element.

Let C (resp. S) be a simply-connected Lie group whose Lie algebra is isomorphic to the centre (resp. the derived algebra) of L(G). Then C is a vector group, S is a compact group with finite centre (Th. 1) and G can be identified with the quotient of C × S by a discrete subgroup D, which is central (Integration, Chap. VII, §3, no. 2, Lemma 4). Since the image of the projection of D onto S is central, hence finite, D∩C is of finite index in D. Let C' be the vector subspace of C generated by D ∩ C, and V a complementary subspace. Then the group T = C'/(D ∩ C) is a torus, and G is isomorphic to the quotient of the product group V × T × S by a finite group.

If G is compact, so is \(V \times T \times S\) (General Topology, Chap. III, §4, no. 1, Cor. 2 of Prop. 2), hence so is V, which implies that \(V = \{e\}\) .

COROLLARY 1. Let G be a connected compact Lie group. Then \(\mathrm{C}(\mathrm{G})_{0}\) is a torus, \(\mathrm{D}(\mathrm{G})\) is a connected compact semi-simple Lie group and the morphism \((x,y)\mapsto xy\) from \(\mathrm{C}(\mathrm{G})_{0}\times\mathrm{D}(\mathrm{G})\) to G is a finite covering.

With the notation in Prop. 4, we have \(\mathrm{V} = \{e\}\) and the subgroups \(f(\mathrm{T})\) and \(f(\mathrm{S})\) of \(\mathrm{G}\) are compact, hence closed. Thus it suffices to show that \(f(\mathrm{T}) = \mathrm{C}(\mathrm{G})_0\) , \(f(\mathrm{S}) = \mathrm{D}(\mathrm{G})\) . Now, \(\mathrm{L}(\mathrm{G}) = \mathrm{L}(f(\mathrm{T})) \times \mathrm{L}(f(\mathrm{S}))\) ; since \(\mathrm{S}\) is semi-simple and \(\mathrm{T}\) is commutative, this implies that \(\mathrm{L}(f(\mathrm{T})) = \mathcal{C}(\mathrm{L}(\mathrm{G})) = \mathrm{L}(\mathrm{C}(\mathrm{G})_0)\) (Chap. III, §9, no. 3, Prop. 8) and \(\mathrm{L}(f(\mathrm{S})) = \mathcal{DL}(\mathrm{G}) = \mathrm{L}(\mathrm{D}(\mathrm{G}))\) (Chap. III, §9, no. 2, Cor. of Prop. 4), hence the stated assertion.

COROLLARY 2. The centre and the fundamental group of a connected compact semi-simple Lie group are finite. Its universal covering is compact.

With the notation in Prop. 4, the groups V and T are reduced to the identity element; thus S is a universal covering of G, and the fundamental group of G is isomorphic to \(\operatorname{Ker} f\) , hence finite. The centre D of G is discrete since G is semi-simple, so D is finite.

COROLLARY 3. The fundamental group of a connected compact Lie group G is a Z-module of finite type, of rank equal to the dimension of C(G).

Indeed, with the notations in Cor. 1, the fundamental group of \(\mathrm{C}(\mathrm{G})_{0}\) is isomorphic to \(Z^{n}\) , with \(n = \dim \mathrm{C}(\mathrm{G})_{0}\) , and the fundamental group of \(\mathrm{D}(\mathrm{G})\) is finite (Cor. 2).

COROLLARY 4. Let G be a connected compact Lie group. The following conditions are equivalent:

\begin{enumerate}
\def\labelenumi{(\roman{enumi})}
\item
  G is semi-simple;
\item
  C(G) is finite;
\item
  \(\pi_{1}(G)\) is finite.
\end{enumerate}

If \(\mathrm{G}\) is simply-connected, it is semi-simple.

This follows from Cor. 1 to 3.

COROLLARY 5. Let G be a connected compact Lie group. Then Int(G) is the identity component of the Lie group Aut(G) (Chap. III, §10, no. 2).

Let \(f \in \mathrm{Aut}(\mathrm{G})_0\) . Then \(f\) induces an automorphism \(f_1\) of \(\mathrm{C}(\mathrm{G})_0\) and an automorphism \(f_2\) of \(\mathrm{D}(\mathrm{G})\) , and we have \(f_1 \in \mathrm{Aut}(\mathrm{C}(\mathrm{G})_0)_0\) , \(f_2 \in \mathrm{Aut}(\mathrm{D}(\mathrm{G}))_0\) . Since \(\mathrm{Aut}(\mathrm{C}(\mathrm{G})_0)\) is discrete (General Topology, Chap. VII, §2, no. 4, Prop. 5), we have \(f_1 = \mathrm{Id}\) ; since \(\mathrm{D}(\mathrm{G})\) is semi-simple, by Chap. III, §10, no. 2, Cor. 2 of Th. 1 there exists an element \(g\) of \(\mathrm{D}(\mathrm{G})\) such that \(f_2(x) = g x g^{-1}\) for all \(x \in \mathrm{D}(\mathrm{G})\) . For all \(x \in \mathrm{C}(\mathrm{G})_0\) , we have \(g x g^{-1} = x = f_1(x)\) ; since \(\mathrm{G} = \mathrm{C}(\mathrm{G})_0.\mathrm{D}(\mathrm{G})\) , it follows that \(g x g^{-1} = f(x)\) for all \(x \in \mathrm{G}\) , so \(f = \operatorname{Int} g\) .

PROPOSITION 5. Let \(\mathbf{G}\) be a Lie group whose Lie algebra is compact.

\begin{enumerate}
\def\labelenumi{\alph{enumi})}
\item
  Assume that G is connected. Then G has a largest compact subgroup K; it is connected. There exists a closed central vector subgroup (no. 2) N of G such that G is the direct product \(N \times K\) .
\item
  Assume that the group of connected components of G is finite. Then:
\end{enumerate}

\begin{enumerate}
\def\labelenumi{(\roman{enumi})}
\item
  Every compact subgroup of \(\mathbf{G}\) is contained in a maximal compact subgroup.
\item
  If \(\mathrm{K}_1\) and \(\mathrm{K}_2\) are two maximal compact subgroups of \(\mathbf{G}\) , there exists \(g \in \mathbf{G}\) such that \(\mathrm{K}_2 = g\mathrm{K}_1g^{-1}\) .
\item
  Let \(\mathrm{K}\) be a maximal compact subgroup of \(\mathrm{G}\) . Then \(\mathrm{K} \cap \mathrm{G}_0\) is equal to \(\mathrm{K}_0\) ; it is the largest compact subgroup of \(\mathrm{G}_0\) .
\item
  There exists a closed central vector subgroup N of \(G_{0}\) , normal in G, such that, for any maximal compact subgroup K of G, \(G_{0}\) is the direct product of \(K_{0}\) by N, and G is the semi-direct product of K by N.
\end{enumerate}

\begin{enumerate}
\def\labelenumi{\alph{enumi})}
\item
  We retain the notations of Prop. 4. The projection of \(\mathrm{Ker}f\) onto V is a finite subgroup of the vector group V, hence is reduced to the identity element. It follows that \(\mathrm{Ker}f\) is contained in \(\mathrm{T} \times \mathrm{S}\) , hence that G is the direct product of the vector group \(\mathrm{N} = f(\mathrm{V})\) and the compact group \(\mathrm{K} = f(\mathrm{T} \times \mathrm{S})\) . Every compact subgroup of G has a projection onto N that is reduced to the identity element, hence is contained in K. This proves \(a\) ).
\item
  Assume now that \(G/G_{0}\) is finite. By a), \(G_{0}\) is the direct product of its largest compact subgroup M and a vector subgroup P; the subgroup M of G is clearly normal. Let n be a vector subspace complement of \(L(M)\) in \(L(G)\) , stable under the adjoint representation of G (no. 1 and no. 3, Prop. 3); this is an ideal of \(L(G)\) and we have \(L(G) = L(M) \times n\) . Let N be the integral subgroup of G with Lie algebra n; by Chap. III, §6, no. 6, Prop. 14, it is normal in G. The projection of \(L(G)\) onto \(L(P)\) with kernel \(L(M)\) induces an isomorphism from n to \(L(P)\) ; it follows that the projection of \(G_{0}\) onto P induces an étale morphism from N to P; since P is simply-connected, this is an isomorphism, and N is a vector group. The morphism \((x, y) \mapsto xy\) from \(M \times N\) to \(G_{0}\) is an injective étale morphism (since \(M \cap N\) is reduced to the identity element), hence an isomorphism. It follows that N is a closed subgroup of G and that the quotient G/N is compact, since \(G_{0}/N\) is compact and \(G/G_{0}\) is finite (General Topology, Chap. III, §4, no. 1, Cor. 2 of Prop. 2).
\end{enumerate}

By Integration, Chap. VII, §3, no. 2, Prop. 3, every compact subgroup of G is contained in a maximal compact subgroup, these are conjugate, and for any maximal compact subgroup K of G, G is the semi-direct product of K by N. Since \(G_{0}\) contains N, it is the semi-direct product of N by \(G_{0} \cap K\) ; it follows that \(G_{0} \cap K\) is connected, hence equal to \(K_{0}\) , since \(\mathrm{K}/(\mathrm{G}_{0} \cap \mathrm{K})\) is isomorphic to \(G/G_{0}\) , hence finite; finally, \(K_{0}\) is clearly the largest compact subgroup of \(G_{0}\) by a).

COROLLARY. If N satisfies the conditions of b) (iv), and if \(K_{1}\) and \(K_{2}\) are two maximal compact subgroups of G, there exists \(n \in N\) such that \(nK_{1}n^{-1} = K_{2}\) .

Indeed, by (ii) there exists an element \(g \in G\) such that \(gK_{1}g^{-1} = K_{2}\) ; by (iv), there exists \(n \in N\) and \(k \in K_{1}\) such that g = nk. The element n then has the required properties.

\section{MAXIMAL TORI OF COMPACT LIE GROUPS}